\documentclass[xcolor=table, aspectratio=169]{beamer}
\usepackage{beamerthemesplit}
\usepackage{wrapfig}
\usetheme{SPbGU}
\usepackage{pdfpages}
\usepackage{amsmath}
\usepackage{amssymb}
\usepackage{cmap}
\usepackage[T2A]{fontenc}
\usepackage[utf8]{inputenc}
\usepackage[english]{babel}
\usepackage{indentfirst}
\usepackage{tikz}
\usetikzlibrary{shapes,arrows,arrows.meta,automata,positioning,quotes,backgrounds,decorations.text,decorations.pathmorphing}
\usepackage{multirow}
\usepackage[noend]{algpseudocode}
\usepackage{algorithm}
\usepackage{algorithmicx}
\usepackage{fancyvrb}
\usepackage[linguistics]{forest}
\usepackage{listings}
\usepackage{multicol}
\usepackage{comment}
\usepackage{xspace}
\usepackage{adjustbox}
\usepackage{makecell}
\usepackage{ stmaryrd }
\usepackage{ulem}
\usepackage{wasysym}
\usepackage{amssymb}% http://ctan.org/pkg/amssymb
\usepackage{pifont}% http://ctan.org/pkg/pifont
\usepackage{animate}
\definecolor{ForestGreen}{RGB}{34,139,34}

\usepackage{xspace}
\usepackage{listings}


\newcommand{\cmark}{\textcolor{green}{\ding{51}}}
\newcommand{\xmark}{\textcolor{red}{\ding{55}}}

\newcommand{\todelete}[1]{{\color{red}{#1}}}

\newcommand{\subdued}[1]{{\color{gray}{#1}}}

\newcommand{\backupbegin}{
   \newcounter{finalframe}
   \setcounter{finalframe}{\value{framenumber}}
}
\newcommand{\backupend}{
   \setcounter{framenumber}{\value{finalframe}}
}

\newcommand{\happyCheck}{\color{green}{\checkmark}}
\newcommand{\timeout}{\color{red}{\clock}}

\newcommand{\makenote}[1]{\hfill \footnotesize{#1}}
\newcommand{\strikeoutnote}[1]{\makenote{\strikethrough{#1}}}
\newcommand{\strikethrough}[1]{\sout{#1}}

\newcommand{\lststrikethrough}[1]{\ttfamily\sout{#1}}

\newcolumntype{A}{>{\hb@xt@\z@\bgroup\hss}r<{\egroup}}
\newcolumntype{B}{>{\hb@xt@\z@\bgroup}l<{\hss\egroup}}

\setbeamertemplate{itemize item}[circle]
\setbeamertemplate{enumerate items}[circle]

\lstdefinelanguage{ocanren}{
keywords=[1]{return, do, where, case, run, conde, fresh, let, match, with, when, class, type,
object, method, of, rec, repeat, until, while, \begin{comment}not,\end{comment} do, done, as, val, inherit,
new, module, sig, deriving, datatype, struct, if, then, else, open, private, virtual, include, success, failure,
true, false, mplus},
keywords=[2]{safe},
keywords=[3]{unsafe},
sensitive=true,
commentstyle=\small\itshape\ttfamily,
keywordstyle=[1]\color{blue},
keywordstyle=[2]\color{ForestGreen},
keywordstyle=[3]\color{violet},
identifierstyle=\ttfamily,
basewidth={0.5em,0.5em},
columns=flexible,
mathescape=true,
escapechar=~,
fontadjust=true,
literate={fun}{{$\lambda$}}1 {function}{function}8 {->}{{$\to$}}3 {<-}{{$\leftarrow$}}3 {===}{{$\equiv$}}1 {=/=}{{$\not\equiv$}}1 {|>}{{$\triangleright$}}3 {\\/}{{$\vee$}}2 {/\\}{{$\wedge$}}2 {^}{{$\uparrow$}}1,
morecomment=[s]{(*}{*)},
 moredelim=**[is][\color{red}]{@!}{@}
}

\tikzstyle{processTree} = [
  ->,
  sibling distance=15em,
  scale=0.6,
  every node/.style = {
    shape=rectangle,
    rounded corners=0.05cm,
    draw,
    align=center,
    minimum size=5mm,
    scale=0.6,},
  %level 1/.style={sibling distance=100em}
  ]


\tikzstyle{program} = [
  draw=black,
  thick,
  rectangle,
  rounded corners=1pt,
  inner sep=5pt,
  inner ysep=5pt
  ]

\tikzstyle{goal} = [
  draw=black,
  rectangle,
  rounded corners=1pt,
  inner ysep=0pt,
  ]

\tikzstyle{input} = [
  draw=none,
  rectangle,
  rounded corners=1pt,
  inner sep=2pt,
  inner ysep=2pt,
  fill=green!10,
  minimum height=5mm
  ]


\tikzstyle{transparent} = [
  draw=none,
  inner ysep=3pt
  ]

\lstset{
language=ocanren
}


\DeclareMathOperator{\Term}{\mathcal{T}}
\DeclareMathOperator{\FlatTerm}{\mathcal{FT}}
\DeclareMathOperator{\Var}{\mathbf{Var}}
\DeclareMathOperator{\Cons}{\mathcal{C}}
\DeclareMathOperator{\Kan}{\mathcal{G}}
\DeclareMathOperator{\Fresh}{\mathbf{Fresh}}
\DeclareMathOperator{\Delay}{\mathbf{Delay}}
\DeclareMathOperator{\Cll}{\mathbf{Call}}
\DeclareMathOperator{\Def}{\mathcal{D}}
\DeclareMathOperator{\Base}{\mathbf{Base}}
\DeclareMathOperator{\Conj}{\mathbf{Conj}}
\DeclareMathOperator{\free}{\mathbf{free}}
\DeclareMathOperator{\ground}{\mathbf{ground}}
\DeclareMathOperator{\In}{\mathbf{In}}
\DeclareMathOperator{\Out}{\mathbf{Out}}
\DeclareMathOperator{\Fun}{\mathcal{F}}
\DeclareMathOperator{\Rtrn}{\mathbf{return}}
\DeclareMathOperator{\Bind}{\mathbf{Bind}}
\DeclareMathOperator{\Match}{\mathbf{Match}}
\DeclareMathOperator{\Sum}{\mathbf{<|>}}
\DeclareMathOperator{\Guard}{\mathbf{Guard}}
\DeclareMathOperator{\Gen}{\mathbf{Gen}}
\DeclareMathOperator{\Stream}{\mathit{Stream}}
\DeclareMathOperator{\vars}{vars}
\DeclareMathOperator{\inmode}{\mathtt{I}}
\DeclareMathOperator{\outmode}{\mathtt{O}}
\DeclareMathOperator{\whatmode}{\mathtt{?}}
% \DeclareMathOperator{\inmode}{g \rightarrow g}
% \DeclareMathOperator{\outmode}{f \rightarrow g}
% \DeclareMathOperator{\mode1}{mode}
% \DeclareMathOperator{\Mode1}{\mathcal{M}}
\newcommand{\KanN}{\mathcal{K}^{N}}
\newcommand{\tran}[1]{\left\llbracket #1 \right\rrbracket}
\newcommand{\LIST}[1]{ #1^*}
\renewcommand{\emptyset}{\varnothing}
\newcommand{\mk}{\textsc{miniKanren}\xspace}
\newcommand{\ocaml}{\textsc{OCaml}\xspace}
\newcommand{\haskell}{\textsc{Haskell}\xspace}
\renewcommand{\and}{$\&$\xspace}
\newcommand{\rel}[2]{\texttt{#1}$^o$ #2}
\newcommand{\subst}[1]{$\langle$#1$\rangle$}
\newcommand{\sem}[1]{\llbracket #1 \rrbracket}

\lstset{moredelim=[is][\bfseries]{[*}{*]}}

\beamertemplatenavigationsymbolsempty

\title[Functional Conversion for miniKanren]{Integration of Offline Partial Deduction and \\ Functional Conversion for miniKanren}
\institute[Jetbrains Research]{
JetBrains Research, Programming Lanuages and Program Analysis Lab
}

\author[Aleksandr Shefer]{Aleksandr Shefer, Ekaterina Verbitskaia}

\date{April 7, 2024}

\definecolor{orange}{RGB}{179,36,31}

\begin{document}
{
\begin{frame}[fragile]
   \begin{center}
      \includegraphics[height=1.5cm]{pictures/jetbrainsResearch.pdf}
    \end{center}
  \titlepage
\end{frame}
}

\definecolor{badRed}{RGB}{225,0,0}
\definecolor{darkFastGreen}{RGB}{10,46,54}

\newcommand{\badBadge}[1]
{
  \begin{tikzpicture}
    \node[circle,
          scale=0.8,
          text width=1cm,
          fill=badRed,
          draw=darkFastGreen,
          line width=0.4mm,
          inner sep=1pt,
          align=center]
        {\textbf{#1}};
  \end{tikzpicture}
}

\definecolor{fastGreen}{RGB}{41,231,204}
\definecolor{darkFastGreen}{RGB}{10,46,54}

\newcommand{\goodBadge}[1]
{
  \begin{tikzpicture}
    \node[circle,
          scale=0.8,
          text width=1cm,
          fill=fastGreen,
          draw=darkFastGreen,
          line width=0.4mm,
          inner sep=1pt,
          align=center]
        {\textbf{#1}};
  \end{tikzpicture}
}

\begin{frame}[fragile]
  \frametitle{Search is More Complicated than Verification}
     % \hfill 
    
    \begin{center}
      \animategraphics[loop,width=4cm]{10}{fig/hanoi/frame_}{0}{26}
    \vfill
         The Tower of Hanoi Puzzle
    \end{center}

    \vfill 

    \lstinline{verify [1->2, 1->3, 2->3] $\Rightarrow$ True }
    
    \lstinline{verify [1->2, 1->2] $\Rightarrow$ False }

    \vfill 
    
    \lstinline{solve True $\Rightarrow$ [1 -> 2, 1->3, 2->3, ...] }
\end{frame}


\begin{frame}[fragile]
  \frametitle{Search is Dual to Verification}
    \begin{center} 
      solve \approx verify$^{-1}$
    \end{center}
\end{frame}

\begin{frame}[fragile]
  \frametitle{Logic Programming Highlights the Duality}

\lstinline{hanoi$^o$ candidate result}

\vfill

\lstinline{verify = run q (hanoi$^o$ [1->2,1->3,2->3] q) $\Rightarrow$ q = True} 

\vfill

\lstinline{solve  = run q (hanoi$^o$ q True) $\Rightarrow$ q = [1 ->2,1->3,2->3,...]}

\end{frame}

% \begin{frame}[fragile]
%     \frametitle{Verifier-to-Solver}

%     
\begin{columns}
    \begin{column}[t]{0.45\textwidth}
        \centering
        Verifier 
        

\begin{figure}[!t]
  \centering
  \begin{minipage}{0.55\textwidth}
    \begin{lstlisting}[frame=tb]
 eval (Conj x y) =
   eval x && eval y
 ...
    \end{lstlisting}
  \end{minipage}
\end{figure}

    \end{column}
    \begin{column}[t]{0.55\textwidth}
        \centering
        Solver 
        
\begin{figure}[!t]
  \centering
  \begin{minipage}{0.7\textwidth}
    \begin{lstlisting}[frame=tb] 
  synth st res = do 
    (u, v) <- [(u, v) | 
      u <- [False, True], 
      v <- [False, True], 
      u && v == res]
    x <- synth st u 
    y <- synth st v 
    return (Conj x y) 
    ...
    \end{lstlisting}
  \end{minipage}
\end{figure}

    \end{column}
\end{columns}
% \end{frame}


% \begin{frame}[fragile]
%     \frametitle{Logic Programming}

%     \begin{columns}[T]
%         \begin{column}{0.49\textwidth}
%             \lstinline{run res (append$^o$ [0] [1] res)} \\
%             \vspace{0.5cm}
%             \lstinline{$\langle$ res $\to$ [0, 1] $\rangle$}
%         \end{column}

%         \begin{column}{0.49\textwidth}
%             \lstinline{run x,y (append$^o$ x y [0, 1])} \\
%             \vspace{0.5cm}
%             \lstinline{$\langle$ x $\to$ [], y $\to$ [0, 1] $\rangle$} \\
%             \lstinline{$\langle$ x $\to$ [0],   y $\to$ [1] $\rangle$} \\
%             \lstinline{$\langle$ x $\to$ [0, 1], y $\to$ [] $\rangle$}
%         \end{column}
%     \end{columns}
% \end{frame}



\begin{frame}[fragile]
    \frametitle{Program Interpretation and Synthesis}

    
\begin{columns}
    \begin{column}[t]{0.45\textwidth}
        \centering
        Verifier 
        

\begin{figure}[!t]
  \centering
  \begin{minipage}{0.55\textwidth}
    \begin{lstlisting}[frame=tb]
 eval (Conj x y) =
   eval x && eval y
 ...
    \end{lstlisting}
  \end{minipage}
\end{figure}

    \end{column}
    \begin{column}[t]{0.55\textwidth}
        \centering
        Solver 
        
\begin{figure}[!t]
  \centering
  \begin{minipage}{0.7\textwidth}
    \begin{lstlisting}[frame=tb] 
  synth st res = do 
    (u, v) <- [(u, v) | 
      u <- [False, True], 
      v <- [False, True], 
      u && v == res]
    x <- synth st u 
    y <- synth st v 
    return (Conj x y) 
    ...
    \end{lstlisting}
  \end{minipage}
\end{figure}

    \end{column}
\end{columns}
\end{frame}

\begin{frame}[fragile]
  \frametitle{Relational Interpreters for Search}

\begin{columns}[T]
  \begin{column}{0.49\textwidth}
    

\begin{figure}[!t]
  \centering
  \begin{minipage}{0.55\textwidth}
    \begin{lstlisting}[frame=tb]
 eval (Conj x y) =
   eval x && eval y
 ...
    \end{lstlisting}
  \end{minipage}
\end{figure}

  \end{column}

  \begin{column}{0.49\textwidth}
    \begin{figure}[!t]
  \centering
  \begin{minipage}{0.75\textwidth}
    \begin{lstlisting}[frame=tb]
 eval$^o$ fm u = fresh (x y v w)
     (fm === Conj x y /\ 
      and$^o$ v w u /\ 
      eval$^o$ x v /\ 
      eval$^o$ y w);
     ...
    \end{lstlisting}
  \end{minipage}
\end{figure}

  \end{column}
\end{columns}

\tikzset{
    block/.style={
      rectangle, draw, minimum width = 2.5cm, minimum height=1.5cm, align=center
        },
}

\newcommand{\fTo}[2]{
  \draw[-Stealth,thick,color=teal] (#1) -- (#2);
}

\begin{center}
  \begin{tikzpicture}[
    every node/.style = {
        inner sep=0pt
      , shape=rectangle
      }
  ]

      \node[block] (int) at (0,0) {miniKanren \\ interpreter};

      \node[] (p_o) at ($(int.west)+(-1cm,0.3cm)$) {eval$^o$};

      \node[] (y) at ($(int.west)+(-1cm,-0.3cm)$) {y};

      \node[] (x) at ($(int.east)+(1cm,0cm)$) {x};
      
      \node[block] (relInt) at ($(p_o.center)+(-2.8,0)$) {relational \\ conversion};

      \node[] (p) at ($(relInt.west)+(-1cm,0cm)$) {eval};

      \fTo{p}{$(p)+(1,0)$};
      \fTo{$(p_o)+(-1.5,0)$}{p_o};
      \fTo{p_o}{$(p_o)+(1,0)$};
      \fTo{y}{$(y)+(1,0)$};
      \fTo{$(x)-(1,0)$}{x};

  \end{tikzpicture}
  \end{center}
  
   
\end{frame}

% \begin{frame}[fragile]
%     \frametitle{Verifier-to-Solver}
    
%     \begin{figure}[h]
  \centering
  \begin{subfigure}[b]{0.45\textwidth}
    \begin{lstlisting}[frame=tb]
let rec eval$^o$ st$^{g \to g}$ fm$^{f \to g}$ u$^{g \to g}$ =
  ( fm$^{f \to g}$  ===    Lit u$^{g \to g}$) \/
  ( elem$^o$ z$^{f \to g}$ st$^{g \to g}$ u$^{g \to g}$ /\
    fm$^{f \to g}$ ===    Var z$^{g \to g}$ ) \/
  ( not$^o$ v$^{f \to g}$ u$^{g \to g}$ /\
    eval$^o$ st$^{g \to g}$ x$^{f \to g}$ v$^{g \to g}$ /\
    fm$^{f \to g}$ ===    Neg x$^{g \to g}$ ) \/
  ( or$^o$ v$^{f \to g}$ w$^{f \to g}$ u$^{g \to g}$ /\
    eval$^o$ st$^{g \to g}$ x$^{f \to g}$ v$^{g \to g}$ /\
    eval$^o$ st$^{g \to g}$ y$^{f \to g}$ w$^{g \to g}$ /\
    fm$^{f \to g}$ ===    Disj x$^{g \to g}$ y$^{g \to g}$) \/
  ( and$^o$ v$^{f \to g}$ w$^{f \to g}$ u$^{g \to g}$ /\
    eval$^o$ st$^{g \to g}$ x$^{f \to g}$ v$^{g \to g}$ /\
    eval$^o$ st$^{g \to g}$ y$^{f \to g}$ w$^{g \to g}$ /\
    fm$^{f \to g}$ ===    Conj x$^{g \to g}$ y$^{g \to g}$) \/
    \end{lstlisting}
   \caption{Mode inference result}
    \label{fig:prop_modded}
  \end{subfigure}
  \hfill
  \begin{subfigure}[b]{0.45\textwidth}
    \begin{lstlisting}[frame=tb]
evaloIOI x0 x2 = msum
  [ do { let {x1 = Lit x2}
       ; return x1 }
  , do { x7 <- elemoOII x0 x2
       ; let {x1 = Var x7}
       ; return x1 }
  , do { x5 <- notoOI x2
       ; x3 <- evaloIOI x0 x5
       ; let {x1 = Neg x3}
       ; return x1 }
  , do { (x5, x6) <- oroOOI x2
       ; x3 <- evaloIOI x0 x5
       ; x4 <- evaloIOI x0 x6
       ; let {x1 = Disj x3 x4}
       ; return x1 }
  , do { (x5, x6) <- andoOOI x2
       ; x3 <- evaloIOI x0 x5
       ; x4 <- evaloIOI x0 x6
       ; let {x1 = Conj x3 x4}
       ; return x1 } ]
    \end{lstlisting}
    \caption{Functional conversion}
    \label{fig:prop_hsk}
  \end{subfigure}
  \caption{Evaluator of propositional formulas}
  \label{fig:prop}
\end{figure}

    
%     \vfill
    
%     \begin{center}
%       \begin{tabular}{c|c} 
%         verifier & solver \\ \hline
%         \lstinline[]$run res (eval$^o$ $[]$ (Disj Trueo Falseo) res)$ & 
%         \lstinline[]$run fm (eval$^o$ $[]$ fm Trueo)$
%       \end{tabular}
%     \end{center}
% \end{frame}


\begin{frame}[fragile]
  \frametitle{Relational Interpreters for Search: the Issue}
\begin{center}
  It is slow
\end{center}
\end{frame}

\begin{frame}[fragile]
    \frametitle{Relational Programming in \mk}

\begin{itemize}
    \item Pure logic programming
    \item Complete search: interleaving
    % \item Easily extensible 
    % \item Embedded
\end{itemize}    
\end{frame}


\begin{frame}[fragile]
    \frametitle{\mk Syntax}

    \begin{center}
\begin{tikzpicture}[
  every node/.style = {
    shape=rectangle,
    rounded corners=0.05cm,
    draw,
    align=center,
    minimum size=5mm},
  node distance=1.3em,
  anchor=center
]
  \node[inner sep=10pt,draw=none] (code) at (0,0) {\begin{lstlisting}
let rec add$^o$ x y z =
  (x === O /\ y === z) \/
  (fresh (x$_1$ z$_1$)
    (x === S x$_1$ /\
     add$^o$ x$_1$ y z$_1$ /\
     z === S z$_1$) )
    \end{lstlisting}};
    \node (rel) [draw=none, above=of code, xshift=1cm] {relation};
    \draw [teal,<-] ($(code.north)+(-0.2,-0.45)$) to [out=90,in=180] node[below, draw=none, teal] {} ($(rel.west)+(0,0)$);
    
    \pause
    \node (call) [draw=none, below=of code, xshift=-2cm] {relation call};
    \draw [teal,<-] ($(code.south)+(-1.2,1.1)$) to [out=-180,in=90] node[below, draw=none, teal] {} ($(call.north)+(0,0)$);
    
    \pause
    \node (uni) [draw=none, left=of code, yshift=1.2cm] {term unification};
    \draw [teal,<-] ($(code.west)+(1.45,0.8)$) to [out=150,in=-45] node[below, draw=none, teal] {} ($(uni.east)+(0,0)$);
    
    \pause
    \node (conj) [draw=none, right=of code, yshift=-0.5cm] {conjunction};
    \draw [teal,<-] ($(code.east)+(-0.8,-0.75)$) to [out=0,in=-180] node[below, draw=none, teal] {} ($(conj.west)+(0,0)$);
    \draw [teal,<-] ($(code.east)+(-1.4,-0.3)$) to [out=0,in=-190] node[below, draw=none, teal] {} ($(conj.west)+(0,0)$);
    
    \pause
    \node (disj) [draw=none, right=of code, yshift=2cm] {disjunction};
    \draw [teal,<-] ($(code.east)+(-0.55,0.85)$) to [out=90,in=-180] node[below, draw=none, teal] {} ($(disj.west)+(0,0)$);
    
    \pause
    \node (fresh) [draw=none, left=of code, yshift=-1cm] {variable introduction};
    \draw [teal,<-] ($(code.west)+(1.1,0)$) to [out=230,in=0] node[below, draw=none, teal] {} ($(fresh.east)+(0,0)$);
\end{tikzpicture}
\end{center}

    
\end{frame}


% \begin{frame}[fragile]
%     \frametitle{Approximation of the Search Strategy}

%     \tikzset{
    block/.style={
      rectangle, minimum width = 2cm, minimum height=1cm, align=center
      %, rounded corners=0.5cm
      %,draw
    },
    block1/.style={
      align=center
    },
}

\newcommand{\fTO}[2]{
  \draw[-Stealth,thick,color=teal] (#1) -- (#2);
}

  % decoration = {
  %   snake,
  %   pre length=2pt,
  %   post length=4pt,
  %   amplitude=0.5pt,
  %   segment length=4pt,
  % },
  
\begin{center}
  \begin{tikzpicture}[
    every node/.style = {
        inner sep=10pt
      , shape=rectangle
      }
  ]

      \coordinate (shift) at (1.5,0);  
      \coordinate (shift1) at (3,0);
      \coordinate (shift2) at (2,0);
      \coordinate (shift3) at (1,0);
      \coordinate (shift4) at (1,0);
      \coordinate (down) at (0,1.8);
      

      \node[block] (main) at ($(0,1)$) {$?\ eval^o([], v_0, True)$};

      \pause


      % $(main)  - (shift1) - (down) - (shift)+(0,0.25)$ 
      % $(main)  - (down) - (shift)$
      % $(main)  - (down) + (shift)$
      % $(main) + (shift1) + (shift) - (down)+(0,0.25)$

      \node[block] (conj) at ($(main)  - (shift1) - (down) - (shift)-(0,0.2)$) {$v_0 \rightarrow Conj(v_1, v_2) $ \\ $eval^o([], v_1, v_3)$ \\ ...};

      \fTO{main.south}{$(conj.north)$};

      
      \node[block] (lit) at ($(main)  - (down) - (shift) +(0,0.2)$) {\underline{$v_0 \rightarrow Lit(True)$}};
      
      \fTO{main.south}{$(lit.north)$};

      
      \node[block] (elem) at ($(main)  - (down) + (shift)$) {$v_0 \rightarrow Var(v_1) $ \\ $elem^o(v_1, [], True)$};


      \fTO{main.south}{$(elem.north)$};

      
      \node[block] (other) at ($(main) + (shift1) + (shift) - (down)+(0,0.25)$) {$...$};

      \fTO{main.south}{$(other.north)$};

      \pause 
      
      \node[block] (conj1) at ($(main)  - (shift1) - (down) - (shift) -(1.0,2.8)$) {$v_0 \rightarrow Conj(v_1, v_2) $ \\ $v_1 \rightarrow Conj(v_5, v_6) $ \\ $eval^o([], v_5, v_7) $ \\ ...};

      \fTO{conj.south}{$(conj1.north)$};
      
      \node[block] (other) at ($(main) - (shift) - (down) -(1.5,2.1)$) {$...$};
      
      \fTO{conj.south}{$(other.north)$};

      \pause 

      \node[block] (fail) at ($(main)  - (down) + (shift)-(0,2.1)$) {$Fail$};

      \fTO{elem.south}{$(fail.north)$};

            
      \pause 

      \node[block] (stop) at ($(main)  - (down) - (down) + (shift)-(0,2.1)$) {\textcolor{red}{\Large When to stop?} \par};


  \end{tikzpicture}
\end{center}
  

%   \begin{tikzpicture}[
%   processTree,
%   sibling distance=10em,
%   level distance=8em]
%   \node {
%     \underline{\rel{eval}{$fm \ s \ true$}}}
%     child { node {
%       \rel{lookup}{$x \ s \ true$} \\
%       \subst{$fm \to var \ v$}}}
%     child { node {
%       \rel{eval}{$x \ s \ a$} \and \\
%       \rel{not}{$a \ true$} \\
%       \subst{$fm \to neg \ x$}}}
%     child { node {
%       \rel{eval}{$x \ s \ a$} \and \\
%       \rel{eval}{$y \ s \ b$} \and \\
%       \rel{or}{a \ b \ true} \\
%       \subst{$fm \to disj \ x \ y$}}}
%     child { node {
%       \rel{eval}{$x \ s \ a$} \and \\
%       \rel{eval}{$y \ s \ b$} \and \\
%       \rel{and}{a \ b \ true} \\
%       \subst{$fm \to conj \ x \ y$}}}
%   ;
% \end{tikzpicture}
% \end{frame}


\begin{frame}[fragile]
    \frametitle{Sources of Inefficiency}

    \begin{itemize}
        \item Running backwards is slow
        \item Order of clauses influences performance
        \item Constant arguments (\lstinline{eval$^{o}$ [] q True)
    \end{itemize}

    \vspace{1cm}

    Solution: specialization
\end{frame}



\begin{frame}[fragile]
  \frametitle{Specialization Scheme}
  \begin{center}
    \tikzset{
    block/.style={
      rectangle, draw, minimum width = 2.5cm, minimum height=1.5cm, align=center
        },
}

\tikzset{
    block1/.style={
      rectangle, draw, minimum width = 2.5cm, minimum height=1.5cm, align=center,color=violet
        },
}

\newcommand{\fTo}[2]{
  \draw[-Stealth,thick,color=teal] (#1) -- (#2);
}


\tikzset{onslide/.code args={<#1>#2}{% from https://tex.stackexchange.com/a/6155/263192
  \only<#1>{\pgfkeysalso{#2}}
}}

\begin{center}
  \begin{tikzpicture}[
    every node/.style = {
        inner sep=0pt
      , shape=rectangle
      }
  ]

      \coordinate (shift) at (2.8,-1.5);  

      \node<2,3>[block] (mode) at ($(relInt)+(shift)$) {mode \\ analysis};
      \node<1>[block1] (mode) at ($(relInt)+(shift)$) {mode \\ analysis};
      
      \node<1>[block1] (annotation) at ($(mode)+(shift)$) {binding time \\ analysis};
      \node<2,3>[block] (annotation) at ($(mode)+(shift)$) {binding time \\ analysis};


      \node<2>[block1] (deduction) at ($(annotation)+(shift)$) {partial \\ deduction};
      \node<1,3>[block] (deduction) at ($(annotation)+(shift)$) {partial \\ deduction};
      \node<1,2>[block] (haskell) at ($(deduction)+(shift)$) {functional \\ conversion to \\ \haskell};
      \node<3>[block1] (haskell) at ($(deduction)+(shift)$) {functional \\ conversion to \\ \haskell};
      
      \node[] (p_o) at ($(mode.west)+(-0.6,0.2)$) {p$^o$};
      \node[] (p_m) at ($(mode.west)!(annotation.north)!(mode.east)$) {p$^m$};
      \node[] (p_a) at ($(annotation.west)!(deduction.north)!(annotation.east)$) {p$^a$};
      \node[] (p_d) at ($(deduction.west)!(haskell.north)!(deduction.east)$) {p$^d$};
      \node[] (p_h) at ($(haskell.east) + (0.6, 0)$) {p$^h$};

      \node[] (dir) at ($(mode.west)+(-1.1,-0.3)$) {direction};
      \node[align=center] (vars) at ($(deduction.west)+(-1.3,0)$) {values of \\ ground vars};

      \fTo{p_o}{$(mode.west) + (0, 0.2)$};

      \fTo{$(mode.east)$}{p_m};
      \fTo{p_m}{annotation.north};

      \fTo{$(annotation.east)$}{p_a};
      \fTo{p_a}{deduction.north};

      \fTo{$(deduction.east)$}{p_d};
      \fTo{p_d}{haskell.north};

      \fTo{dir}{$(mode.west) + (0.0,-0.3)$};
      \fTo{vars}{$(deduction.west)$};
      
      \fTo{$(haskell.east)$}{p_h};

      
      \pause 
      % \node[block] (mode) at ($(relInt)+(shift)$) {mode \\ analysis};
      % \node[block] (annotation) at ($(mode)+(shift)$) {binding time \\ analysis};

      \pause 
      % \node[block] (deduction) at ($(annotation)+(shift)$) {partial \\ deduction};
      

  \end{tikzpicture}
  \end{center}
  
  

  \end{center}
\end{frame}

\begin{frame}[fragile]
  \frametitle{Mode Analysis}
  \begin{center}
    \begin{minipage}{0.5\textwidth}
      \begin{columns}[T]
        \begin{column}{0.1\textwidth}
          Variable modes: 
        \end{column}
        \begin{column}{0.6\textwidth}
         \begin{minipage}{\textwidth}
            \begin{itemize}
                \item \underline{Ground} term: True
                \item \underline{Free} variable: x
            \end{itemize}
          \end{minipage}
        \end{column}
    \end{columns}
  \end{minipage}
  \end{center}

\vspace{0.3cm}
  
  \begin{center}
    \lstinline{eval$^o$ [] fm True -> eval$^o$ g f g}
    
    \begin{center}
  \begin{tikzpicture}[
    every node/.style = {
        inner sep=0pt
      , shape=rectangle
      % , draw
      }
  ]
      % \draw[help lines] (-7.75,-3.5) grid (7.75,3.5);
      % \node[inner sep=10pt, anchor=north west] (func) at (-7.75,3.5) {\begin{minipage}{0.35\columnwidth}
  \begin{lstlisting}[frame=tb]
 let rec mult$^o$ x y z = conde [
   (fresh (x$_1$ r$_1$)
     (x === S x$_1$) /\
     (add$^o$ y r$_1$ z) /\
     (mult$^o$ x$_1$ y r$_1$));
   ...]
\end{lstlisting}
\end{minipage}};

      % \node[inner sep=10pt, anchor=north east] (slow) at (7.75,3.5) {\begin{minipage}{0.52\columnwidth}
  \begin{lstlisting}[frame=tb]
 multIIO$_1$ $::$ Nat -> Nat -> Stream Nat
 multIIO$_1$ (S x$_1$) y    = do
   (r$_1$, r) <- addIOO y
   multIII x$_1$ y r$_1$
   return r
 ...
 multIII $::$ Nat -> Nat -> Nat -> Stream ()
 multIII (S x$_1$) y z = do
   z$_1$ <- multIIO$_1$ x$_1$ y
   addIII y z$_1$ z
 multIII _ _ _ = Empty
 ...
\end{lstlisting}
\end{minipage}
};

      \node[inner sep=10pt] (main) at (0,0)
      {\begin{columns}
  \centering

  \begin{column}{0.35\textwidth}
    \begin{lstlisting}[frame=tb]
   eval$^o$ st fm u = 
     fresh (x y v w z) 
       ([*fm === Conj x y*] /\
        eval$^o$ st x v /\
        eval$^o$ st y w /\ 
        [*and$^o$ v w u*]) \/
        ... 
    \end{lstlisting}
  \end{column}
  
  \begin{column}{0.35\textwidth}
    \begin{lstlisting}[frame=tb]
   eval$^o$ st fm u = 
     fresh (x y v w z) 
       ([*and$^o$ v w u*] /\
        eval$^o$ st x v /\
        eval$^o$ st y w /\ 
        [*fm === Conj x y*]) \/
        ... 
    \end{lstlisting}
  \end{column}
\end{columns}};
      % \node [draw=none, above=of main.south, xshift=2.3cm] {$\Omega(x!)$};

      \node[ellipse, scale=0.9,
          fill=fastGreen,
          line width=0.4mm,
          inner sep=1pt,
          align=center, above=of main.south, xshift=7.0cm, yshift=1.84cm]
        {\lstinline{u = True}};

        
      \node[ellipse, scale=0.9,
          fill=fastGreen,
          line width=0.4mm,
          inner sep=1pt,
          align=center, above=of main.south, xshift=7.0cm, yshift=1.26cm]
        {\lstinline{v = True}};
        
      \node[ellipse, scale=0.9,
          fill=fastGreen,
          line width=0.4mm,
          inner sep=1pt,
          align=center, above=of main.south, xshift=7.0cm, yshift=0.68cm]
        {\lstinline{w = True}};

        
      \node[ellipse, scale=0.9,
          fill=badRed!75,
          line width=0.4mm,
          inner sep=1pt,
          align=center, above=of main.south, xshift=-1.0cm, yshift=0.26cm]
        {\lstinline{u = True}};

      % \node[inner sep=10pt, anchor=south west] (fast) at (-7.75,-3.5) {\begin{minipage}{0.42\columnwidth}
\begin{lstlisting}[frame=tb]
multIIO $::$ Nat -> Nat -> Stream Nat
multIIO (S x$_1$) y = do
  r$_1$ <- multIIO x$_1$ y
  addIIO y r$_1$
...
\end{lstlisting}
\end{minipage}
};

      % \node [draw=none, anchor=south west] at ($(main.south)+(0.5,0.8)$) {\goodBadge{$O(xy)$}};
      % \node [draw=none, above=of main.south, xshift=2.3cm] {\goodBadge{10x \linebreak faster}};
  \end{tikzpicture}
  \end{center}
  
  \end{center}

\end{frame}

\begin{frame}[fragile]
  \frametitle{Partial Deduction}

  \begin{tikzpicture}[
  node distance = 14mm and 13 mm,
  decoration = {
    snake,
    pre length=2pt,
    post length=4pt,
    amplitude=0.5pt,
    segment length=4pt
  },
  remember picture,overlay]
  \node (a) [
    program,
    anchor=north west,
    xshift=2cm,
    yshift=-1.8cm]
  at (current page.north west)
  {
    \adjustbox{scale=0.65}{
      \begin{minipage}[c]{0.38\textwidth}
        \begin{lstlisting}
    eval$^o$ st fm u =
      safe and$^o$ a b u &
      unsafe eval$^o$ st x a &
      unsafe eval$^o$ st y b &
      fm === Conj x y |
      ...
        \end{lstlisting}
      \end{minipage}
    }
  };

  \node (b) [goal,anchor=north east] at (a.south east)
  {
    \adjustbox{scale=0.65}{
      \begin{minipage}[c]{0.24\textwidth}
        \begin{lstlisting}
    eval$^o$ [] fm true
        \end{lstlisting}
      \end{minipage}
    }
  };

  \pause

  \node (d) [
    transparent,
    xshift=-2cm,
    yshift=-2.8cm,
    anchor=north east]
  at (current page.north east)
  {
      \begin{tikzpicture}[
  processTree,
  sibling distance=7em,
  level distance=7em,
  level 2/.style={level distance=5em},
  anchor=center]
  \node {
    \rel{eval}{$[]\ fm \ true$}}
    child { node[draw=none, fill=none] {...}}
    child { node {
      \rel{and}{$a \ b \ true$} \and  \\
      \rel{eval}{$[]\ x \ a$} \and \\
      \rel{eval}{$[]\ y \ b$} \\
      \subst{$fm \to Conj \ x \ y$}}
      child { node[draw=none, fill=none] {...}}
      }
    child { node[draw=none, fill=none] {...}}
  ;
\end{tikzpicture}
  };

  \draw[->,semithick, decorate]
    (a.north east) to
    [out=15,in=145,"{\footnotesize constructing trees}"]
    ($(d.north west)+(0.9,0)$);

  % \pause

  % \node (e) [transparent, anchor=north east] at ($(d.south east)+(0,-1)$)
  % {
  %   \begin{minipage}[c]{0.4\textwidth}
  %     \begin{tikzpicture}[
  processTree,
  sibling distance=4em,
  level distance=3em,
  level 3/.style={level distance=4em,sibling distance=8em},
  anchor=center]
  \node (root) {
    \rel{eval}{$fm \ s \ true$}}
    child { node[draw=none, fill=none] {...}}
    child { node[draw=none, fill=none] {...}
      child { node {
        \rel{eval}{$x \ s \ true$} \and \\
        \rel{eval}{$y \ s \ true$}}
        child { node (1) {\rel{eval}{$x \ s \ true$}}}
        child { node (2) {\rel{eval}{$y \ s \ true$}}}
        }
    }
    child { node[draw=none, fill=none] {...}}
  ;

  \draw [dashed,->] (1.west) to [out=170,in=-150] (root.west);
  \draw [dashed,->] (2.east) to [out=10,in=-30] (root.east);


\end{tikzpicture}
  %   \end{minipage}
  % };

  % \draw[->,semithick, decorate]
  %   ($(d.south)+(0.6,0.7)$) to
  %   [out=-75,in=45,"{\footnotesize folding}"]
  %   (e.north);

  \pause


  \node (c) [
    program,
    anchor=south west,
    xshift=2cm,
    yshift=0.8cm]
  at (current page.south west)
  {
    \adjustbox{scale=0.65}{
      \begin{minipage}[c]{0.4\textwidth}
        \begin{lstlisting}
    eval$^o$_true fm  =
      eval$^o$_true x &
      eval$^o$_true y & 
      fm === conj x y |
      ...
        \end{lstlisting}
      \end{minipage}
    }
  };

  \draw[->,semithick, decorate]
    ($(d.south)+(-1.0,0)$) to
    [out=230,in=-10,"{\footnotesize residualization}",pos=0.4]
    (c.east);

    \onslide<1->
\end{tikzpicture}
\end{frame}


\begin{frame}[fragile]
  \frametitle{Partial Deduction: Trees Construction}

  \tikzset{
    block/.style={
      rectangle, minimum width = 2cm, minimum height=1cm, align=center
      %, rounded corners=0.5cm
      %,draw
    },
    block1/.style={
      rectangle, minimum width = 1cm, minimum height=1cm, align=center
    },
}

\newcommand{\fTO}[2]{
  \draw[-Stealth,thick,color=teal] (#1) -- (#2);
}

\newcommand{\fTORed}[2]{
  \draw[-Stealth,thick,color=violet] (#1) -- (#2);
}

\newcommand{\fTOGreen}[2]{
  \draw[-Stealth,thick,color=ForestGreen] (#1) -- (#2);
}

  % decoration = {
  %   snake,
  %   pre length=2pt,
  %   post length=4pt,
  %   amplitude=0.5pt,
  %   segment length=4pt,
  % },
  
\begin{center}
  \begin{tikzpicture}[
    every node/.style = {
        inner sep=10pt
      , shape=rectangle
      }
  ]

      \coordinate (shift) at (1.7,0);  
      \coordinate (shift1) at (3.4,0);
      \coordinate (shift2) at (2,0);
      \coordinate (shift3) at (1,0);
      \coordinate (shift4) at (1,0);
      \coordinate (down) at (0,1.8);
      

      \node[block] (main) at ($(0,1)$) {$?\ eval^o([], v_0, True)$};


      % $(main)  - (shift1) - (down) - (shift)+(0,0.25)$ 
      % $(main)  - (down) - (shift)$
      % $(main)  - (down) + (shift)$
      % $(main) + (shift1) + (shift) - (down)+(0,0.25)$

      \node[block] (conj) at ($(main)  - (shift1) - (down) - (shift)-(0,0.5)$) { $and^o(v_3, v_5, True)$ \and \\ $eval^o([], v_1, v_3)$ \and \\ $eval^o([], v_2, v_4)$ \\ \subst{$v_0 \rightarrow Conj(v_1, v_2) $}};

      \fTO{main.south}{$(conj.north)$};

      
      \node[block] (lit) at ($(main)  - (down) - (shift) +(0,0.2)$) {\subst{\underline{$v_0 \rightarrow Lit(True)$}}};
      
      \fTO{main.south}{$(lit.north)$};

      
      \node[block] (elem) at ($(main)  - (down) + (shift)$) {$elem^o(v_1, [], True)$ \\ \subst{$v_0 \rightarrow Var(v_1) $}};


      \fTO{main.south}{$(elem.north)$};

      
      \node[block] (other) at ($(main) + (shift1) + (shift) - (down)+(0,0.25)$) {$...$};

      \fTO{main.south}{$(other.north)$};
      
      
      \node[block] (fail) at ($(main)  - (down) + (shift)-(0,2.1)$) {$Fail$};

      \fTO{elem.south}{$(fail.north)$};

      \node[block1] (stop) at ($(main)- (shift1) - (down) - (down) - (shift)$) {\textcolor{red}{Stop!}};
       
  \end{tikzpicture}
\end{center}
  

%   \begin{tikzpicture}[
%   processTree,
%   sibling distance=10em,
%   level distance=8em]
%   \node {
%     \underline{\rel{eval}{$fm \ s \ true$}}}
%     child { node {
%       \rel{lookup}{$x \ s \ true$} \\
%       \subst{$fm \to var \ v$}}}
%     child { node {
%       \rel{eval}{$x \ s \ a$} \and \\
%       \rel{not}{$a \ true$} \\
%       \subst{$fm \to neg \ x$}}}
%     child { node {
%       \rel{eval}{$x \ s \ a$} \and \\
%       \rel{eval}{$y \ s \ b$} \and \\
%       \rel{or}{a \ b \ true} \\
%       \subst{$fm \to disj \ x \ y$}}}
%     child { node {
%       \rel{eval}{$x \ s \ a$} \and \\
%       \rel{eval}{$y \ s \ b$} \and \\
%       \rel{and}{a \ b \ true} \\
%       \subst{$fm \to conj \ x \ y$}}}
%   ;
% \end{tikzpicture}
\end{frame}

\begin{frame}[fragile]
  \frametitle{Binding Time Analysis}

  \begin{center}
    \begin{minipage}{0.5\textwidth}
      \begin{columns}[T]
        \begin{column}{0.3\textwidth}
          Call 
          
          annotations: 
        \end{column}
        \begin{column}{0.8\textwidth}
         \begin{minipage}{\textwidth}
            \begin{itemize}
                \item \underline{Safe}: continue unfolding
                \item \underline{Unsafe}: \textbf{possibly} stop unfolding 
            \end{itemize}
          \end{minipage}
        \end{column}
    \end{columns}
  \end{minipage}
  \end{center}

  \vspace{0.3cm}

  \begin{center}
    \begin{minipage}[c]{0.38\textwidth}
        \begin{lstlisting}[frame=tb]
   eval$^o$ st fm u =
    fresh (x y v w z)
     (safe and$^o$ a b u &
      unsafe eval$^o$ st x a &
      unsafe eval$^o$ st y b &
      fm === Conj x y) |
      ...
        \end{lstlisting}
      \end{minipage}
  \end{center}  

\end{frame}

\begin{frame}[fragile]
  \frametitle{Partial Deduction with Annotations}

  \tikzset{
    block/.style={
      rectangle, minimum width = 2cm, minimum height=1cm, align=center
      %, rounded corners=0.5cm
      %,draw
    },
    block1/.style={
      rectangle, minimum width = 1cm, minimum height=1cm, align=center
    },
}

\newcommand{\fTO}[2]{
  \draw[-Stealth,thick,color=teal] (#1) -- (#2);
}

\newcommand{\fTORed}[2]{
  \draw[-Stealth,thick,color=violet] (#1) -- (#2);
}

\newcommand{\fTOGreen}[2]{
  \draw[-Stealth,thick,color=ForestGreen] (#1) -- (#2);
}

  % decoration = {
  %   snake,
  %   pre length=2pt,
  %   post length=4pt,
  %   amplitude=0.5pt,
  %   segment length=4pt,
  % },
  
\begin{center}
  \begin{tikzpicture}[
    every node/.style = {
        inner sep=10pt
      , shape=rectangle
      }
  ]

      \coordinate (shift) at (1.7,0);  
      \coordinate (shift1) at (3.4,0);
      \coordinate (shift2) at (2,0);
      \coordinate (shift3) at (1,0);
      \coordinate (shift4) at (1,0);
      \coordinate (down) at (0,1.8);
      

      \node[block] (main) at ($(0,1)$) {$?\ eval^o([], v_0, True)$};


      % $(main)  - (shift1) - (down) - (shift)+(0,0.25)$ 
      % $(main)  - (down) - (shift)$
      % $(main)  - (down) + (shift)$
      % $(main) + (shift1) + (shift) - (down)+(0,0.25)$

      \node[block] (conj) at ($(main)  - (shift1) - (down) - (shift)-(0,0.5)$) { $and^o(v_3, v_5, True)$ \and \\ $eval^o([], v_1, v_3)$ \and \\ $eval^o([], v_2, v_4)$ \\ \subst{$v_0 \rightarrow Conj(v_1, v_2) $}};

      \fTO{main.south}{$(conj.north)$};

      
      \node[block] (lit) at ($(main)  - (down) - (shift) +(0,0.2)$) {\subst{\underline{$v_0 \rightarrow Lit(True)$}}};
      
      \fTO{main.south}{$(lit.north)$};

      
      \node[block] (elem) at ($(main)  - (down) + (shift)$) {$elem^o(v_1, [], True)$ \\ \subst{$v_0 \rightarrow Var(v_1) $}};


      \fTO{main.south}{$(elem.north)$};

      
      \node[block] (other) at ($(main) + (shift1) + (shift) - (down)+(0,0.25)$) {$...$};

      \fTO{main.south}{$(other.north)$};
      
      
      \node[block] (fail) at ($(main)  - (down) + (shift)-(0,2.1)$) {$Fail$};

      \fTO{elem.south}{$(fail.north)$};

      
      \node[block] (conj1) at ($(main)  - (shift1) - (down) - (down) - (0,1.0) - (shift)-(0,0.5)$) {$eval^o([], v_1, True)$ \and \\ $eval^o([], v_2, True)$ \\ \subst{$v_0 \rightarrow Conj(v_1, v_2) $}};

      \fTO{conj.south}{$(conj1.north)$};

      \node[block1] (unsafe) at ($(main) - (down) - (down) - (0,1.0) - (shift)-(0,0.5)$) {\textcolor{violet}{Unsafe}};

      \fTORed{unsafe.west}{$(conj1.east)$};

      
      \node[block1] (unsafe) at ($(main) - (down) - (shift) - (0,0.5)$) {\textcolor{ForestGreen}{Safe}};

      \fTOGreen{unsafe.west}{$(conj.east)$};


      \node[block1] (stop) at ($(main)- (shift1) - (down) - (down) - (down) - (shift) - (0,0.8)$) {\textcolor{red}{Stop!}};
       
  \end{tikzpicture}
\end{center}
  

%   \begin{tikzpicture}[
%   processTree,
%   sibling distance=10em,
%   level distance=8em]
%   \node {
%     \underline{\rel{eval}{$fm \ s \ true$}}}
%     child { node {
%       \rel{lookup}{$x \ s \ true$} \\
%       \subst{$fm \to var \ v$}}}
%     child { node {
%       \rel{eval}{$x \ s \ a$} \and \\
%       \rel{not}{$a \ true$} \\
%       \subst{$fm \to neg \ x$}}}
%     child { node {
%       \rel{eval}{$x \ s \ a$} \and \\
%       \rel{eval}{$y \ s \ b$} \and \\
%       \rel{or}{a \ b \ true} \\
%       \subst{$fm \to disj \ x \ y$}}}
%     child { node {
%       \rel{eval}{$x \ s \ a$} \and \\
%       \rel{eval}{$y \ s \ b$} \and \\
%       \rel{and}{a \ b \ true} \\
%       \subst{$fm \to conj \ x \ y$}}}
%   ;
% \end{tikzpicture}
\end{frame}


\begin{frame}[fragile]
  \frametitle{Functional Conversion}

  \begin{columns}[T]
  \centering
  
  \begin{column}{0.35\textwidth}
    \begin{lstlisting} 
  eval_true$^o$ fm = 
    fresh (x y) 
      (eval_true$^o$ x /\
       eval_true$^o$ y /\ 
       fm === Conj x y) \/
      ... 
    \end{lstlisting}
  \end{column}

  \begin{column}{0.1\textwidth}
    \begin{center}
      $\Rightarrow$
    \end{center}
  \end{column}

  \begin{column}{0.35\textwidth}
    \begin{lstlisting}
  eval_true = do 
    x <- eval_true
    y <- eval_true 
    return (Conj x y) 
    ... 
    \end{lstlisting}
  \end{column}
\end{columns}  
\end{frame}

\begin{frame}[fragile]
  \frametitle{Evaluation}
  \begin{itemize}
      \item Functional Conversion
      \item Online Partial Deduction and Functional Conversion
      \item \underline{Our approach}: Offline Partial Deduction and Functional Conversion
  \end{itemize}
  
\end{frame}


\begin{frame}[fragile]
  \frametitle{Evaluation: Propositional Evaluator}
  \begin{center}
    \begin{table}[]
\begin{tabular}{r||r||r|r}

& \multicolumn{1}{c||}{\multirow[t]{2}{*}{\begin{tabular}[c]{@{}l@{}}\texttt{Functional}\\ \text{conversion}\end{tabular}}} & \multicolumn{2}{c}{\begin{tabular}[c]{@{}l@{}}\texttt{Translation with}\\ \text{partial deduction}\end{tabular}} \\ \cline{3-4} & 
        \multicolumn{1}{c||}{} & \multicolumn{1}{c|}{Online} & \multicolumn{1}{c}{Offline} \\ \hline
      

                % & \multicolumn{2}{c||}{\begin{tabular}[c]{@{}l@{}}\texttt{Translation with}\\ \text{partial deduction}\end{tabular}} & \multicolumn{1}{c}{\multirow[t]{2}{*}{\begin{tabular}[c]{@{}l@{}}\texttt{Functional}\\ \text{conversion}\end{tabular}}} \\ \cline{2-3}
                % & \multicolumn{1}{r|}{Online} & \multicolumn{1}{r||}{Offline} & \multicolumn{1}{c}{}                            \\ \hline
\texttt{10 formulas}  & \SI{4.05}{\micro\second} & \multicolumn{1}{r|}{\SI{0.33}{\micro\second}} & \SI{0.37}{\micro\second} \\
\texttt{100 formulas} & \SI{56.00}{\micro\second} & \multicolumn{1}{r|}{\SI{7.40}{\micro\second}}  & \SI{7.73}{\micro\second}  \\
\texttt{1000 formulas} & \SI{645.00}{\micro\second} & \multicolumn{1}{r|}{\SI{108.00}{\micro\second}}  & \SI{108.00}{\micro\second}  
\end{tabular}
\end{table}
  \end{center}
  
\end{frame}

\begin{frame}[fragile]
  \frametitle{Evalutaion: the Tower of Hanoi puzzle}
  \begin{center}
    \begin{table}[]
  \begin{tabular}{l||r||r|r}
        & \multicolumn{1}{c||}{\multirow[t]{2}{*}{\begin{tabular}[c]{@{}l@{}}\texttt{Functional}\\ \text{conversion}\end{tabular}}} & \multicolumn{2}{c}{\begin{tabular}[c]{@{}l@{}}\texttt{Translation with}\\ \text{partial deduction}\end{tabular}} \\ \cline{3-4} & 
        \multicolumn{1}{c||}{} & \multicolumn{1}{c|}{Online} & \multicolumn{1}{c}{Offline} \\ \hline
      \texttt{3 disks} & \SI{153000.00}{\micro\second} & \SI{125000.00}{\micro\second} & \SI{1.67}{\micro\second} \\
      \texttt{4 disks} & \timeout timeout  & \timeout timeout & \SI{3.12}{\micro\second} \\
  \end{tabular}
\end{table}
  \end{center}
\end{frame}

\begin{frame}[fragile]
  \frametitle{Evalutaion: Other Benchmarks (DPPD\footnote{Dozens of Problems for Partial Deduction benchmark: \url{https://github.com/leuschel/DPPD}})}
  \begin{center}
    \begin{table}[]
  \begin{tabular}{l||r||r|r}
        & \multicolumn{1}{c||}{\multirow[t]{2}{*}{\begin{tabular}[c]{@{}l@{}}\texttt{Functional}\\ \text{conversion}\end{tabular}}} & \multicolumn{2}{c}{\begin{tabular}[c]{@{}l@{}}\texttt{Translation with}\\ \text{partial deduction}\end{tabular}} \\ \cline{3-4} & 
        \multicolumn{1}{c||}{} & \multicolumn{1}{c|}{Online} & \multicolumn{1}{c}{Offline} \\ \hline
      \texttt{appLast} & \SI{1.13}{\micro\second} & \SI{0.08}{\micro\second} & \SI{0.09}{\micro\second} \\
      \texttt{contains} & \SI{14.70}{\micro\second}  & \SI{0.78}{\micro\second} & \SI{0.77}{\micro\second} \\
      \texttt{doubleAppend} & \timeout timeout  & \SI{2910.00}{\micro\second} & \SI{2.00}{\micro\second} \\
      \texttt{exDepth} & \timeout timeout  & \timeout timeout & \SI{3.09}{\micro\second} \\
      \texttt{nthOpt} & \SI{0.10}{\micro\second}  & \SI{0.08}{\micro\second} & \SI{0.10}{\micro\second} \\
      ... &  &  &   \\

  \end{tabular}
\end{table}
  \end{center}
\end{frame}


\begin{frame}[fragile]
  \frametitle{Conclusion}

  \begin{itemize}
      \item We integrated the Functional Conversion and Offline Partial Deduction approaches
      \item We conducted a preliminary evaluation of the approach
      \item The effectiveness of the integration has been shown
  \end{itemize}

  
\end{frame}

% \begin{frame}[fragile]
%     \frametitle{Prior Research}
    
%     \begin{itemize}
%         \item Michael Leuschel, Stephen J. Craig, Maurice Bruynooghe, and Wim Vanhoof. 2004. Specialising Interpreters using Offline Partial Deduction
%         \item Ekaterina Verbitskaia, Daniil Berezun, and Dmitry Boulytchev. 2021. An Empirical Study of Partial Deduction for miniKanren
%         \item Ekaterina Verbitskaia, Igor Engel, and Daniil Berezun. 2024. A Case Study in Functional Conversion and Mode Inference in miniKanren

%     \end{itemize}
  
% \end{frame}


% \begin{frame}[fragile]
%     \frametitle{Goal}

%     \textbf{Goal}: Develop a more powerful specialization technique for relational programming.

%     \vfill
    
%     \textbf{Objectives}: 
%     \begin{itemize}
%         \item Analyze common specialization techniques in logic and relational programming. 
%         \item Tune the most suitable approaches to suit our goal best. 
%         \item Implement a transformation method which combines specialization and functional conversion.
%         \item Collect benchmark suite. Evaluate the proposed approach. 
%     \end{itemize}
% \end{frame}

% \begin{frame}[fragile]
%   \frametitle{Algorithm scheme}
%   \begin{center}
%     \tikzset{
    block/.style={
      rectangle, draw, minimum width = 2.5cm, minimum height=1.5cm, align=center
        },
}

\tikzset{
    block1/.style={
      rectangle, draw, minimum width = 2.5cm, minimum height=1.5cm, align=center,color=violet
        },
}

\newcommand{\fTo}[2]{
  \draw[-Stealth,thick,color=teal] (#1) -- (#2);
}


\tikzset{onslide/.code args={<#1>#2}{% from https://tex.stackexchange.com/a/6155/263192
  \only<#1>{\pgfkeysalso{#2}}
}}

\begin{center}
  \begin{tikzpicture}[
    every node/.style = {
        inner sep=0pt
      , shape=rectangle
      }
  ]

      \coordinate (shift) at (2.8,-1.5);  

      \node<2,3>[block] (mode) at ($(relInt)+(shift)$) {mode \\ analysis};
      \node<1>[block1] (mode) at ($(relInt)+(shift)$) {mode \\ analysis};
      
      \node<1>[block1] (annotation) at ($(mode)+(shift)$) {binding time \\ analysis};
      \node<2,3>[block] (annotation) at ($(mode)+(shift)$) {binding time \\ analysis};


      \node<2>[block1] (deduction) at ($(annotation)+(shift)$) {partial \\ deduction};
      \node<1,3>[block] (deduction) at ($(annotation)+(shift)$) {partial \\ deduction};
      \node<1,2>[block] (haskell) at ($(deduction)+(shift)$) {functional \\ conversion to \\ \haskell};
      \node<3>[block1] (haskell) at ($(deduction)+(shift)$) {functional \\ conversion to \\ \haskell};
      
      \node[] (p_o) at ($(mode.west)+(-0.6,0.2)$) {p$^o$};
      \node[] (p_m) at ($(mode.west)!(annotation.north)!(mode.east)$) {p$^m$};
      \node[] (p_a) at ($(annotation.west)!(deduction.north)!(annotation.east)$) {p$^a$};
      \node[] (p_d) at ($(deduction.west)!(haskell.north)!(deduction.east)$) {p$^d$};
      \node[] (p_h) at ($(haskell.east) + (0.6, 0)$) {p$^h$};

      \node[] (dir) at ($(mode.west)+(-1.1,-0.3)$) {direction};
      \node[align=center] (vars) at ($(deduction.west)+(-1.3,0)$) {values of \\ ground vars};

      \fTo{p_o}{$(mode.west) + (0, 0.2)$};

      \fTo{$(mode.east)$}{p_m};
      \fTo{p_m}{annotation.north};

      \fTo{$(annotation.east)$}{p_a};
      \fTo{p_a}{deduction.north};

      \fTo{$(deduction.east)$}{p_d};
      \fTo{p_d}{haskell.north};

      \fTo{dir}{$(mode.west) + (0.0,-0.3)$};
      \fTo{vars}{$(deduction.west)$};
      
      \fTo{$(haskell.east)$}{p_h};

      
      \pause 
      % \node[block] (mode) at ($(relInt)+(shift)$) {mode \\ analysis};
      % \node[block] (annotation) at ($(mode)+(shift)$) {binding time \\ analysis};

      \pause 
      % \node[block] (deduction) at ($(annotation)+(shift)$) {partial \\ deduction};
      

  \end{tikzpicture}
  \end{center}
  
  

%   \end{center}
% \end{frame}


% \begin{frame}[fragile]
%   \frametitle{Evalutaion: Propositional evaluator}
%   \begin{center}
%     \begin{table}[]
\begin{tabular}{r||r||r|r}

& \multicolumn{1}{c||}{\multirow[t]{2}{*}{\begin{tabular}[c]{@{}l@{}}\texttt{Functional}\\ \text{conversion}\end{tabular}}} & \multicolumn{2}{c}{\begin{tabular}[c]{@{}l@{}}\texttt{Translation with}\\ \text{partial deduction}\end{tabular}} \\ \cline{3-4} & 
        \multicolumn{1}{c||}{} & \multicolumn{1}{c|}{Online} & \multicolumn{1}{c}{Offline} \\ \hline
      

                % & \multicolumn{2}{c||}{\begin{tabular}[c]{@{}l@{}}\texttt{Translation with}\\ \text{partial deduction}\end{tabular}} & \multicolumn{1}{c}{\multirow[t]{2}{*}{\begin{tabular}[c]{@{}l@{}}\texttt{Functional}\\ \text{conversion}\end{tabular}}} \\ \cline{2-3}
                % & \multicolumn{1}{r|}{Online} & \multicolumn{1}{r||}{Offline} & \multicolumn{1}{c}{}                            \\ \hline
\texttt{10 formulas}  & \SI{4.05}{\micro\second} & \multicolumn{1}{r|}{\SI{0.33}{\micro\second}} & \SI{0.37}{\micro\second} \\
\texttt{100 formulas} & \SI{56.00}{\micro\second} & \multicolumn{1}{r|}{\SI{7.40}{\micro\second}}  & \SI{7.73}{\micro\second}  \\
\texttt{1000 formulas} & \SI{645.00}{\micro\second} & \multicolumn{1}{r|}{\SI{108.00}{\micro\second}}  & \SI{108.00}{\micro\second}  
\end{tabular}
\end{table}
%   \end{center}
% \end{frame}


% \begin{frame}[fragile]
%   \frametitle{Evalutaion: Tower of Hanoi puzzle}
%   \begin{center}
%     \begin{table}[]
  \begin{tabular}{l||r||r|r}
        & \multicolumn{1}{c||}{\multirow[t]{2}{*}{\begin{tabular}[c]{@{}l@{}}\texttt{Functional}\\ \text{conversion}\end{tabular}}} & \multicolumn{2}{c}{\begin{tabular}[c]{@{}l@{}}\texttt{Translation with}\\ \text{partial deduction}\end{tabular}} \\ \cline{3-4} & 
        \multicolumn{1}{c||}{} & \multicolumn{1}{c|}{Online} & \multicolumn{1}{c}{Offline} \\ \hline
      \texttt{3 disks} & \SI{153000.00}{\micro\second} & \SI{125000.00}{\micro\second} & \SI{1.67}{\micro\second} \\
      \texttt{4 disks} & \timeout timeout  & \timeout timeout & \SI{3.12}{\micro\second} \\
  \end{tabular}
\end{table}
%   \end{center}
% \end{frame}

% \begin{frame}[fragile]
%   \frametitle{Evalutaion: other benchmarks (DPPD)}
%   \begin{center}
%     \begin{table}[]
  \begin{tabular}{l||r||r|r}
        & \multicolumn{1}{c||}{\multirow[t]{2}{*}{\begin{tabular}[c]{@{}l@{}}\texttt{Functional}\\ \text{conversion}\end{tabular}}} & \multicolumn{2}{c}{\begin{tabular}[c]{@{}l@{}}\texttt{Translation with}\\ \text{partial deduction}\end{tabular}} \\ \cline{3-4} & 
        \multicolumn{1}{c||}{} & \multicolumn{1}{c|}{Online} & \multicolumn{1}{c}{Offline} \\ \hline
      \texttt{appLast} & \SI{1.13}{\micro\second} & \SI{0.08}{\micro\second} & \SI{0.09}{\micro\second} \\
      \texttt{contains} & \SI{14.70}{\micro\second}  & \SI{0.78}{\micro\second} & \SI{0.77}{\micro\second} \\
      \texttt{doubleAppend} & \timeout timeout  & \SI{2910.00}{\micro\second} & \SI{2.00}{\micro\second} \\
      \texttt{exDepth} & \timeout timeout  & \timeout timeout & \SI{3.09}{\micro\second} \\
      \texttt{nthOpt} & \SI{0.10}{\micro\second}  & \SI{0.08}{\micro\second} & \SI{0.10}{\micro\second} \\
      ... &  &  &   \\

  \end{tabular}
\end{table}
%   \end{center}
% \end{frame}

% \begin{frame}[fragile]
%   \frametitle{Results (6 months)}
%   \begin{itemize}
%       \item Surveyed and adapted the relevant techniques for my task
%       \item Collected benchmark suite
%       \item Conducted a preliminary evaluation
%   \end{itemize}

%   \vfill

%   Future work: 
%   \begin{itemize}
%       \item Improve the quality of the mode analysis stage 
%       \item Finish evaluation 
%   \end{itemize}
% \end{frame}


% \begin{frame}[fragile]
%   \frametitle{Inverse Interpreter}
%   \begin{columns}[t]
%     \begin{column}{0.49\textwidth}
%       \begin{center} 
%         \setlength{\tabcolsep}{0.1cm}
%         \setlength{\extrarowheight}{10pt}
%         \begin{tabular}{rcl}
%           $\llbracket p \rrbracket (x)$ & $=$ & $y$ \\ 
%           $\llbracket p^{-1} \rrbracket (y)$ & $=$ & $x$ \\ 
%           $\llbracket invInt \rrbracket (p, y)$ & $=$ & $x$
%         \end{tabular}
%       \end{center}
%     \end{column}
%     \begin{column}{0.49\textwidth}
%       \vspace{0.5cm}
%       \tikzset{
    block/.style={
      rectangle, draw, minimum width = 2.5cm, minimum height=1.5cm, align=center
        },
}

\newcommand{\fTo}[2]{
  \draw[-Stealth,thick,color=teal] (#1) -- (#2);
}

\begin{center}
  \begin{tikzpicture}[
    every node/.style = {
        inner sep=0pt
      , shape=rectangle
      }
  ]
      \node[block] (int) at (0,0) {Inverse \\ interpreter};

      \node[] (p) at ($(int.west)+(-1cm,0.3cm)$) {p};
      \node[] (y) at ($(int.west)+(-1cm,-0.3cm)$) {y};

      \node[] (x) at ($(int.east)+(1cm,0cm)$) {x};
     
      \fTo{p}{$(p)+(1,0)$};
      \fTo{y}{$(y)+(1,0)$};
      \fTo{$(x)-(1,0)$}{x};

  \end{tikzpicture}
  \end{center}


\begin{center}
  \begin{tikzpicture}[
    every node/.style = {
        inner sep=0pt
      , shape=rectangle
      }
  ]
      \node[block] (int) at (0,0) {miniKanren};

      \node[] (p) at ($(int.west)+(-1cm,0.3cm)$) {p};
      \node[] (y) at ($(int.west)+(-1cm,-0.3cm)$) {y};

      \node[] (x) at ($(int.east)+(1cm,0cm)$) {x};
     
      \fTo{p}{$(p)+(1,0)$};
      \fTo{y}{$(y)+(1,0)$};
      \fTo{$(x)-(1,0)$}{x};

  \end{tikzpicture}
  \end{center}
  
  
%     \end{column}
%   \end{columns}
% \end{frame}

% \begin{frame}[fragile]
%   \frametitle{miniKanren as an Inverse Interpreter}
%   \begin{columns}[t]
%     \begin{column}{0.49\textwidth}
%       \begin{figure}[!t]
  \centering
  \begin{minipage}{0.7\columnwidth}
    \begin{lstlisting}[frame=tb]
 let rec add$^o$ x y z = conde [
   (x === O /\ y === z);
   (fresh (x$_1$ z$_1$)
     (x === S x$_1$ /\
      add$^o$ x$_1$ y z$_1$ /\
      z === S z$_1$) ) ]
    \end{lstlisting}
  \end{minipage}
\end{figure}
%    \end{column}
%     \begin{column}{0.49\textwidth}
%       \lstinline{run z (add$^o$ 0 1 z) = $\{$z $\mapsto$ 1$\}$}

%       \tikzset{
    block/.style={
      rectangle, draw, minimum width = 2.5cm, minimum height=1.5cm, align=center
        },
}

\newcommand{\fTo}[2]{
  \draw[-Stealth,thick,color=teal] (#1) -- (#2);
}

\begin{center}
  \begin{tikzpicture}[
    every node/.style = {
        inner sep=0pt
      , shape=rectangle
      }
  ]
      \node[block] (int) at (0,0) {miniKanren \\ interpreter};

      \node[] (p) at ($(int.west)+(-1cm,0.5cm)$) {add$^o$};
      \node[] (x) at ($(int.west)+(-1cm,0cm)$) {x};
      \node[] (y) at ($(int.west)+(-1cm,-0.5cm)$) {y};

      \node[] (z) at ($(int.east)+(1cm,0cm)$) {z};
     
      \fTo{p}{$(p)+(1,0)$};
      \fTo{x}{$(x)+(1,0)$};
      \fTo{y}{$(y)+(1,0)$};
      \fTo{$(z)-(1,0)$}{z};

  \end{tikzpicture}
  \end{center}
  
  

%       \vspace{0.5cm}

    
%       \lstinline{run x, y (add$^o$ x y 1) =} 
\lstinline{       [$\{$x $\mapsto$ 0, y $\mapsto$ 1$\}$, $\{$x $\mapsto$ 1, y $\mapsto$ 0$\}$]}

%       \tikzset{
    block/.style={
      rectangle, draw, minimum width = 2.5cm, minimum height=1.5cm, align=center
        },
}

\begin{center}
  \begin{tikzpicture}[
    every node/.style = {
        inner sep=0pt
      , shape=rectangle
      }
  ]
      \node[block] (int) at (0,0) {miniKanren \\ interpreter};

      \node[] (p) at ($(int.west)+(-1cm,0.3cm)$) {add$^o$};
      \node[] (z) at ($(int.west)+(-1cm,-0.3cm)$) {z};

      \node[] (xy) at ($(int.east)+(1cm,0cm)$) {(x, y)};
     
      \fTo{p}{$(p)+(1,0)$};
      \fTo{z}{$(z)+(1,0)$};
      \fTo{$(xy)-(1,0)$}{xy};

  \end{tikzpicture}
  \end{center}
  
  
%     \end{column}
%   \end{columns}
% \end{frame}

% \begin{frame}[fragile]
%   \frametitle{Relational Interpreters for Search}

% \begin{columns}
%   \begin{column}{0.49\textwidth}
%     

\begin{figure}[!t]
  \centering
  \begin{minipage}{0.55\textwidth}
    \begin{lstlisting}[frame=tb]
 eval (Conj x y) =
   eval x && eval y
 ...
    \end{lstlisting}
  \end{minipage}
\end{figure}

%   \end{column}

%   \begin{column}{0.49\textwidth}
%     \begin{figure}[!t]
  \centering
  \begin{minipage}{0.75\textwidth}
    \begin{lstlisting}[frame=tb]
 eval$^o$ fm u = fresh (x y v w)
     (fm === Conj x y /\ 
      and$^o$ v w u /\ 
      eval$^o$ x v /\ 
      eval$^o$ y w);
     ...
    \end{lstlisting}
  \end{minipage}
\end{figure}

%   \end{column}
% \end{columns}

% \tikzset{
    block/.style={
      rectangle, draw, minimum width = 2.5cm, minimum height=1.5cm, align=center
        },
}

\newcommand{\fTo}[2]{
  \draw[-Stealth,thick,color=teal] (#1) -- (#2);
}

\begin{center}
  \begin{tikzpicture}[
    every node/.style = {
        inner sep=0pt
      , shape=rectangle
      }
  ]

      \node[block] (int) at (0,0) {miniKanren \\ interpreter};

      \node[] (p_o) at ($(int.west)+(-1cm,0.3cm)$) {eval$^o$};

      \node[] (y) at ($(int.west)+(-1cm,-0.3cm)$) {y};

      \node[] (x) at ($(int.east)+(1cm,0cm)$) {x};
      
      \node[block] (relInt) at ($(p_o.center)+(-2.8,0)$) {relational \\ conversion};

      \node[] (p) at ($(relInt.west)+(-1cm,0cm)$) {eval};

      \fTo{p}{$(p)+(1,0)$};
      \fTo{$(p_o)+(-1.5,0)$}{p_o};
      \fTo{p_o}{$(p_o)+(1,0)$};
      \fTo{y}{$(y)+(1,0)$};
      \fTo{$(x)-(1,0)$}{x};

  \end{tikzpicture}
  \end{center}
  
   
% \end{frame}

% \begin{frame}[fragile]
%   \frametitle{Relational Interpreters for Search: the Issue}
% \begin{center}
%   It is slow
% \end{center}

% \vspace{0.5cm}

% % \begin{center}
% %   \begin{minipage}{0.4\textwidth}
% %     \begin{itemize}
% %       \item Unifications 
% %       \item Occurs-check
% %       \item Scheduling complexity
% %     \end{itemize}
% %   \end{minipage}
% % \end{center}
% \end{frame}

% \begin{frame}[fragile]
%   \frametitle{Functional Conversion to the Rescue}
% \begin{center}
%   \tikzset{
    block/.style={
      rectangle, draw, minimum width = 2.5cm, minimum height=1.5cm, align=center
        },
}

\newcommand{\fTo}[2]{
  \draw[-Stealth,thick,color=teal] (#1) -- (#2);
}

\begin{center}
  \begin{tikzpicture}[
    every node/.style = {
        inner sep=0pt
      , shape=rectangle
      }
  ]

      \node[block] (int) at (0,0) {functional \\ conversion};

      \node[] (p_o) at ($(int.west)+(-1.2cm,0.3cm)$) {eval$^o$};

      \node[] (dir) at ($(int.west)+(-1.2cm,-0.3cm)$) {direction};

      \node[] (eval1) at ($(int.east)+(1.3cm,0cm)$) {eval$^{-1}$};
      
      \node[block] (relInt) at ($(p_o.center)+(-2.8,0)$) {relational \\ conversion};

      \node[] (p) at ($(relInt.west)+(-1cm,0cm)$) {eval};

      \fTo{p}{$(p)+(1,0)$};
      \fTo{$(p_o)+(-1.5,0)$}{p_o};
      \fTo{p_o}{$(p_o)+(1.2,0)$};
      \fTo{dir}{$(dir)+(1.2,0)$};
      \fTo{$(eval1)-(1.2,0)$}{eval1};

  \end{tikzpicture}
  \end{center}
  
  
% \end{center}

% \vfill

% \begin{center}
%   \begin{minipage}{0.4\textwidth}
%     \begin{itemize}
%       \item Generate the same answers as \mk would
%       \item Hopefully faster
%     \end{itemize}
%   \end{minipage}
% \end{center}
% \end{frame}

% \begin{frame}[fragile]
%   \frametitle{\mk Syntax}
% \begin{center}
%   \begin{center}
\begin{tikzpicture}[
  every node/.style = {
    shape=rectangle,
    rounded corners=0.05cm,
    draw,
    align=center,
    minimum size=5mm},
  node distance=1.3em,
  anchor=center
]
  \node[inner sep=10pt,draw=none] (code) at (0,0) {\begin{lstlisting}
let rec add$^o$ x y z =
  (x === O /\ y === z) \/
  (fresh (x$_1$ z$_1$)
    (x === S x$_1$ /\
     add$^o$ x$_1$ y z$_1$ /\
     z === S z$_1$) )
    \end{lstlisting}};
    \node (rel) [draw=none, above=of code, xshift=1cm] {relation};
    \draw [teal,<-] ($(code.north)+(-0.2,-0.45)$) to [out=90,in=180] node[below, draw=none, teal] {} ($(rel.west)+(0,0)$);
    
    \pause
    \node (call) [draw=none, below=of code, xshift=-2cm] {relation call};
    \draw [teal,<-] ($(code.south)+(-1.2,1.1)$) to [out=-180,in=90] node[below, draw=none, teal] {} ($(call.north)+(0,0)$);
    
    \pause
    \node (uni) [draw=none, left=of code, yshift=1.2cm] {term unification};
    \draw [teal,<-] ($(code.west)+(1.45,0.8)$) to [out=150,in=-45] node[below, draw=none, teal] {} ($(uni.east)+(0,0)$);
    
    \pause
    \node (conj) [draw=none, right=of code, yshift=-0.5cm] {conjunction};
    \draw [teal,<-] ($(code.east)+(-0.8,-0.75)$) to [out=0,in=-180] node[below, draw=none, teal] {} ($(conj.west)+(0,0)$);
    \draw [teal,<-] ($(code.east)+(-1.4,-0.3)$) to [out=0,in=-190] node[below, draw=none, teal] {} ($(conj.west)+(0,0)$);
    
    \pause
    \node (disj) [draw=none, right=of code, yshift=2cm] {disjunction};
    \draw [teal,<-] ($(code.east)+(-0.55,0.85)$) to [out=90,in=-180] node[below, draw=none, teal] {} ($(disj.west)+(0,0)$);
    
    \pause
    \node (fresh) [draw=none, left=of code, yshift=-1cm] {variable introduction};
    \draw [teal,<-] ($(code.west)+(1.1,0)$) to [out=230,in=0] node[below, draw=none, teal] {} ($(fresh.east)+(0,0)$);
\end{tikzpicture}
\end{center}

% \end{center}
% \end{frame}

% \begin{frame}[fragile]
%   \frametitle{Example: Addition in the Forward Direction}
% \begin{columns}
%   \begin{column}[t]{0.49\textwidth}
%     \begin{figure}[!t]
  \centering
  \begin{minipage}{0.7\columnwidth}
    \begin{lstlisting}[frame=tb]
 let rec add$^o$ x y z = conde [
   (x === O /\ y === z);
   (fresh (x$_1$ z$_1$)
     (x === S x$_1$ /\
      add$^o$ x$_1$ y z$_1$ /\
      z === S z$_1$) ) ]
    \end{lstlisting}
  \end{minipage}
\end{figure}
%     \[ \text{add}^{\circ}\ 0\ 1\ z = \{z \mapsto 1\}\]
%   \end{column}
%   \begin{column}[t]{0.49\textwidth}
%     \begin{figure}[!t]
  \centering
  \begin{minipage}{\columnwidth}
    \begin{lstlisting}[frame=tb]
addIIO :: Nat -> Nat -> Nat
addIIO x y =
  case x of
    O -> y
    S x$_1$ -> S (addIIO x$_1$ y)
    \end{lstlisting}
  \end{minipage}
\end{figure}

%     \[ \text{addIIO}\ 0\ 1 = 1 \] 
%   \end{column}
% \end{columns}
% \end{frame}

% \begin{frame}[fragile]
%   \frametitle{Noninjective $\leftrightarrow$ Nondeterministic}
% \tikzset{
    domainnode/.style={
        shape=ellipse,
        draw, 
        minimum width=1.5cm, 
        minimum height=3cm
        },
}

\newcommand{\point}[3]{
  \node[] (point_#1_#3) at ($(#1.center)+(0,#2)$) {#3};
}

\newcommand{\fTo}[2]{
  \draw[-Stealth,thick,color=teal] (#1) -- (#2);
}

\begin{center}
  \begin{tikzpicture}[
    every node/.style = {
        inner sep=0pt
      , shape=rectangle
      }
  ]

      % \draw[help lines] (-7.5,-3.5) grid (7.5,3.5);
      \node[domainnode] (from) at (-5,0) {};
      \point{from}{1}{1};
      \point{from}{0}{2};
      \point{from}{-1}{3};

      \node[domainnode] (to) at (-2,0) {};
      \point{to}{1}{a};
      \point{to}{0}{b};
      \point{to}{-1}{c};

      \fTo{point_from_1}{point_to_b};
      \fTo{point_from_2}{point_to_c};
      \fTo{point_from_3}{point_to_c};

      \node[] (p) at ($(from.east)!0.5!(to.west) + (0,1)$) {$p$};

      \node[domainnode] (from1) at (2,0) {};
      \point{from1}{1}{1};
      \point{from1}{0}{2};
      \point{from1}{-1}{3};

      \node[domainnode] (to1) at (5,0) {};
      \point{to1}{1}{a};
      \point{to1}{0}{b};
      \point{to1}{-1}{c};

      \fTo{point_to1_b}{point_from1_1};
      \fTo{point_to1_c}{point_from1_2};
      \fTo{point_to1_c}{point_from1_3};

      \node[] (p1) at ($(from1.east)!0.5!(to1.west) + (0,1)$) {$p^{-1}$};

  \end{tikzpicture}
  \end{center}
  

% \vfill

% \begin{center}
%     \[ \text{add}^{\circ}\ x\ y\ 1 = \left[\{x \mapsto  0, y \mapsto 1\}, \{x \mapsto  1, y \mapsto  0\} \right] \] 
% \end{center}
% \end{frame}

% \begin{frame}[fragile]
%   \frametitle{Stream}

%   \begin{center}
%     \begin{minipage}{0.4\textwidth}
%       \begin{itemize}
%           \item Represents nondeterminism
%           \item List-like
%           \item Interleaving search
%       \end{itemize}
%     \end{minipage}
%   \end{center}

% \vfill

%   \begin{center}
%       \[ \left[1, 2, 3\right] >>= f = f(1) < \mid > f(2) < \mid > f(3)\] 
%       \[ \left[1, 2, 3\right] < \mid > \left[a, b, c\right] = \left[1, a, 2, b, 3, c\right] \] 
%   \end{center}

%   \vfill

%   \begin{center}
%     \begin{minipage}{0.5\textwidth}
%       \begin{itemize}
%           \item \mk: \lstinline{Stream} of substitutions
%           \item Functional conversion: \lstinline{Stream} of values
%       \end{itemize}
%     \end{minipage}
%   \end{center}
% \end{frame}


% \begin{frame}[fragile]
%   \frametitle{Addition in the Backward Direction: Nondeterminism}
% \begin{columns}
%   \begin{column}[t]{0.49\textwidth}
%     \begin{figure}[!t]
  \centering
  \begin{minipage}{0.7\columnwidth}
    \begin{lstlisting}[frame=tb]
 let rec add$^o$ x y z = conde [
   (x === O /\ y === z);
   (fresh (x$_1$ z$_1$)
     (x === S x$_1$ /\
      add$^o$ x$_1$ y z$_1$ /\
      z === S z$_1$) ) ]
    \end{lstlisting}
  \end{minipage}
\end{figure}
%   \end{column}
%   \begin{column}[t]{0.49\textwidth}
%     \begin{figure}[!t]
  \centering
  \begin{minipage}{0.9\columnwidth}
    \begin{lstlisting}[frame=tb,language=ocanren1]
addOOI $::$ Nat -> Stream (Nat, Nat)
addOOI z =
  return (O, z) <$\mid$>
  case z of
    O -> Empty
    S z$_1$ -> do
      (x$_1$, y) <- addOOI z$_1$
      return (S x$_1$, y)
    \end{lstlisting}
  \end{minipage}
\end{figure}

%   \end{column}
% \end{columns}

% \vfill
% $\text{add}^{\circ}\ x\ y\ 1 = \left[\{x \mapsto  0, y \mapsto  1\}, \{x \mapsto  1, y \mapsto  0\} \right]$

% \vfill

% \lstinline{addOOI 1 = [(0,1), (1,0)]}

% \end{frame}



% % \begin{frame}[fragile]
% %     \frametitle{Program inversion}
% % \goodBadge{10}
% % \badBadge{$\Omega(n!)$}
% % \end{frame}
% % \begin{frame}[fragile]
% %     \frametitle{Program inversion}
% %     \begin{center}
% %         Many complicated programgs are inverse of simple ones
% %     \end{center}
% %     \vfill
% %     \begin{center}
% %         Type inference or habitation is inverse of type checking
% %     \end{center}
% %     \vfill
% %     \begin{center}
% %         Program inversion: Given a program $f$, produce inverse porgram  $f^{-1}$
% %     \end{center}
% %     \vfill
% %     \begin{center}
% %         Given $\text{typecheck}(\text{program}, \text{types}) = true$, produce $\text{typecheck}^{-1}(\text{program}, true) = \text{types}$.
% %     \end{center}
% % \end{frame}
% % \begin{frame}[fragile]
% %     \frametitle{Relational inversion}
% %     \begin{center}
% %         A program is a relation between its inputs and outputs
% %     \end{center}
% %     \vfill
% %     \begin{center}
% %         miniKanren can evaluate relations in both directions
% %     \end{center}
% %     \vfill
% %     \begin{center}
% %         Write a simple verifier, convert to miniKanren, get a solver
% %     \end{center}
% %     \vfill
% %     \begin{center}
% %         Problem: miniKanren may be slow. So convert it back!
% %     \end{center}
% % \end{frame}

% % \begin{frame}[fragile]
% %   \frametitle{Relational Programming}

% %     \begin{center}
% %         Subset of logic programming, focus on pure relations
% %     \end{center}
% %     \vfill
% %     \begin{center}
% %         Extra-logical features (cuts, side-effects, search order manipulation) discouraged
% %     \end{center}
% %     \vfill
% %     \begin{center}
% %         Interleaving search guarantees that all answers are found
% %     \end{center}
% % \end{frame}

% % \begin{frame}[fragile]
% %   \frametitle{miniKanren}
% %     \begin{center}
% %         miniKanren is a simple relational language designed to be implemented as shallow embedding.
% %     \end{center}
% %     \begin{center}
\begin{tikzpicture}[
  every node/.style = {
    shape=rectangle,
    rounded corners=0.05cm,
    draw,
    align=center,
    minimum size=5mm},
  node distance=1.3em,
  anchor=center
]
  \node[inner sep=10pt,draw=none] (code) at (0,0) {\begin{lstlisting}
let rec add$^o$ x y z =
  (x === O /\ y === z) \/
  (fresh (x$_1$ z$_1$)
    (x === S x$_1$ /\
     add$^o$ x$_1$ y z$_1$ /\
     z === S z$_1$) )
    \end{lstlisting}};
    \node (rel) [draw=none, above=of code, xshift=1cm] {relation};
    \draw [teal,<-] ($(code.north)+(-0.2,-0.45)$) to [out=90,in=180] node[below, draw=none, teal] {} ($(rel.west)+(0,0)$);
    
    \pause
    \node (call) [draw=none, below=of code, xshift=-2cm] {relation call};
    \draw [teal,<-] ($(code.south)+(-1.2,1.1)$) to [out=-180,in=90] node[below, draw=none, teal] {} ($(call.north)+(0,0)$);
    
    \pause
    \node (uni) [draw=none, left=of code, yshift=1.2cm] {term unification};
    \draw [teal,<-] ($(code.west)+(1.45,0.8)$) to [out=150,in=-45] node[below, draw=none, teal] {} ($(uni.east)+(0,0)$);
    
    \pause
    \node (conj) [draw=none, right=of code, yshift=-0.5cm] {conjunction};
    \draw [teal,<-] ($(code.east)+(-0.8,-0.75)$) to [out=0,in=-180] node[below, draw=none, teal] {} ($(conj.west)+(0,0)$);
    \draw [teal,<-] ($(code.east)+(-1.4,-0.3)$) to [out=0,in=-190] node[below, draw=none, teal] {} ($(conj.west)+(0,0)$);
    
    \pause
    \node (disj) [draw=none, right=of code, yshift=2cm] {disjunction};
    \draw [teal,<-] ($(code.east)+(-0.55,0.85)$) to [out=90,in=-180] node[below, draw=none, teal] {} ($(disj.west)+(0,0)$);
    
    \pause
    \node (fresh) [draw=none, left=of code, yshift=-1cm] {variable introduction};
    \draw [teal,<-] ($(code.west)+(1.1,0)$) to [out=230,in=0] node[below, draw=none, teal] {} ($(fresh.east)+(0,0)$);
\end{tikzpicture}
\end{center}

% % \end{frame}


% % \begin{frame}[fragile]
% %   \frametitle{Functional Conversion}
% % \begin{center}
% %   Given a relation and a principal direction, construct a functional program that generates the same answers as \mk would
% % \end{center}

% % \vfill

% % \begin{center}
% %   Preserve the completeness of the search
% % \end{center}

% % \vfill

% % \begin{center}
% % Both inputs and outputs are expected to be ground
% % \end{center}

% % \vfill

% % \begin{center}
% %     Speed improvement: up to $3$ orders of magnitude on benchmark of mulltiplication
% % \end{center}

% % \end{frame}

% \begin{frame}[fragile]
%   \frametitle{Free Variables in Answers: Generators}
% \begin{columns}
%   \begin{column}[t]{0.49\textwidth}
%     \begin{figure}[!t]
  \centering
  \begin{minipage}{0.7\columnwidth}
    \begin{lstlisting}[frame=tb]
 let rec add$^o$ x y z = conde [
   (x === O /\ y === z);
   (fresh (x$_1$ z$_1$)
     (x === S x$_1$ /\
      add$^o$ x$_1$ y z$_1$ /\
      z === S z$_1$) ) ]
    \end{lstlisting}
  \end{minipage}
\end{figure}

%     \vspace{0.5cm}
%     $\text{add}^{\circ}\ 1\ y\ z = \{z \mapsto S\ y\}$

%     \vspace{0.5cm}

%     \lstinline{genNat = [0, 1, 2,$\ldots$]}

%     \vspace{0.5cm}

%     \lstinline{addIOO 1 = [(0,1), (1,2), (2,3),$\ldots$]} 
%   \end{column}
%   \begin{column}[t]{0.49\textwidth}
%     \begin{figure}[!t]
  \centering
  \begin{minipage}{\columnwidth}
    \begin{lstlisting}[label={add_x},
                      %  caption={Function for \lstinline{addo in out out} direction},
                       captionpos=b,
                       frame=tb]
addX :: Nat -> Stream (Nat, Nat)
addX x = case x of
           O -> do
             z <- genNat
             return (z, z)
           S x' -> do
             (y, z') <- addX x'
             return (y, S z')

genNat :: Stream Nat
genNat = Mature O (S <$\$$> genNat)
    \end{lstlisting}
  \end{minipage}
\end{figure}
%   \end{column}
% \end{columns}
% \end{frame}




% \lstset{basicstyle=\small}

% % \begin{frame}[fragile]
% %   \frametitle{Predicates}
% % \begin{columns}
% %   \begin{column}[t]{0.48\textwidth}
% %     \begin{figure}[!t]
  \centering
  \begin{minipage}{0.7\columnwidth}
    \begin{lstlisting}[frame=tb]
 let rec add$^o$ x y z = conde [
   (x === O /\ y === z);
   (fresh (x$_1$ z$_1$)
     (x === S x$_1$ /\
      add$^o$ x$_1$ y z$_1$ /\
      z === S z$_1$) ) ]
    \end{lstlisting}
  \end{minipage}
\end{figure}
% %   \end{column}
% %   \begin{column}[t]{0.5\textwidth}
% %     \begin{figure}[!t]
  \centering
  \begin{minipage}{\columnwidth}
    \begin{lstlisting}[frame=tb]
 addIII $::$ Nat -> Nat -> Nat -> Stream ()
 addIII x y z =
   case x of
     O | y == z -> return ()
       | otherwise -> Empty
     S x$_1$ ->
       case z of
         O -> Empty
         S z$_1$ -> addIII x$_1$ y z$_1$
    \end{lstlisting}
  \end{minipage}
\end{figure}
% %   \end{column}
% % \end{columns}
% % \end{frame}

% \begin{frame}[fragile]
%   \frametitle{Conversion Scheme}
%   \begin{center}
%     \tikzset{
    block/.style={
      rectangle, draw, minimum width = 2.5cm, minimum height=1.5cm, align=center
        },
}

\tikzset{
    block1/.style={
      rectangle, draw, minimum width = 2.5cm, minimum height=1.5cm, align=center,color=violet
        },
}

\newcommand{\fTo}[2]{
  \draw[-Stealth,thick,color=teal] (#1) -- (#2);
}


\tikzset{onslide/.code args={<#1>#2}{% from https://tex.stackexchange.com/a/6155/263192
  \only<#1>{\pgfkeysalso{#2}}
}}

\begin{center}
  \begin{tikzpicture}[
    every node/.style = {
        inner sep=0pt
      , shape=rectangle
      }
  ]

      \coordinate (shift) at (2.8,-1.5);  

      \node<2,3>[block] (mode) at ($(relInt)+(shift)$) {mode \\ analysis};
      \node<1>[block1] (mode) at ($(relInt)+(shift)$) {mode \\ analysis};
      
      \node<1>[block1] (annotation) at ($(mode)+(shift)$) {binding time \\ analysis};
      \node<2,3>[block] (annotation) at ($(mode)+(shift)$) {binding time \\ analysis};


      \node<2>[block1] (deduction) at ($(annotation)+(shift)$) {partial \\ deduction};
      \node<1,3>[block] (deduction) at ($(annotation)+(shift)$) {partial \\ deduction};
      \node<1,2>[block] (haskell) at ($(deduction)+(shift)$) {functional \\ conversion to \\ \haskell};
      \node<3>[block1] (haskell) at ($(deduction)+(shift)$) {functional \\ conversion to \\ \haskell};
      
      \node[] (p_o) at ($(mode.west)+(-0.6,0.2)$) {p$^o$};
      \node[] (p_m) at ($(mode.west)!(annotation.north)!(mode.east)$) {p$^m$};
      \node[] (p_a) at ($(annotation.west)!(deduction.north)!(annotation.east)$) {p$^a$};
      \node[] (p_d) at ($(deduction.west)!(haskell.north)!(deduction.east)$) {p$^d$};
      \node[] (p_h) at ($(haskell.east) + (0.6, 0)$) {p$^h$};

      \node[] (dir) at ($(mode.west)+(-1.1,-0.3)$) {direction};
      \node[align=center] (vars) at ($(deduction.west)+(-1.3,0)$) {values of \\ ground vars};

      \fTo{p_o}{$(mode.west) + (0, 0.2)$};

      \fTo{$(mode.east)$}{p_m};
      \fTo{p_m}{annotation.north};

      \fTo{$(annotation.east)$}{p_a};
      \fTo{p_a}{deduction.north};

      \fTo{$(deduction.east)$}{p_d};
      \fTo{p_d}{haskell.north};

      \fTo{dir}{$(mode.west) + (0.0,-0.3)$};
      \fTo{vars}{$(deduction.west)$};
      
      \fTo{$(haskell.east)$}{p_h};

      
      \pause 
      % \node[block] (mode) at ($(relInt)+(shift)$) {mode \\ analysis};
      % \node[block] (annotation) at ($(mode)+(shift)$) {binding time \\ analysis};

      \pause 
      % \node[block] (deduction) at ($(annotation)+(shift)$) {partial \\ deduction};
      

  \end{tikzpicture}
  \end{center}
  
  


%     % \begin{minipage}{0.4\textwidth}
%     %   \begin{enumerate}
%     %       \item Normalization
%     %       \item Mode analysis
%     %       \item Functional conversion
%     %   \end{enumerate}
%     % \end{minipage}
%   \end{center}
% \end{frame}


% \begin{frame}[fragile]
%   \frametitle{Normalization: Flat Term}
% \begin{center}

% Eliminate nested constructors and repeated variables

% \end{center}

%   \[  \FlatTerm = V \cup \{\Cons \, x_0 \ldots \, x_{k} \mid x_{j}\in V, x_j - distinct \} \]

% \vfill

% \begin{center}
% \begin{tabular}{rcl}
%  $C\left( x, y \right) \equiv C\left( C\left( v, u \right), w \right)$ & $\iff$ &  $x \equiv C\left(v, u \right) \land y \equiv w$ \\
%  $add^{\circ}(x, x, z)$ & $\iff$ & $add^{\circ}\left( x, y, z \right) \land x \equiv y$
% \end{tabular}
% \end{center}

% \end{frame}

% \begin{frame}[fragile]
%   \frametitle{Normalization: Goal}
% \begin{center}
%   Eliminate disjunctions within conjunctions
% \end{center}

% \vfill

% \begin{center}
% \setlength{\tabcolsep}{12pt}
% \begin{tabular}{rclll}
% $\KanN$ & $::=$ & $c_1 \vee \ldots \vee c_{n}$ & $c_{i}\in \Conj$ & \subdued{normal form} \\
% $\Conj$ & $::=$ & $g_1 \wedge \ldots \wedge g_n$ & $ g_{i}\in \Base$ & \subdued{normal conjunction} \\
% $\Base$ & $::=$ & $V \equiv \FlatTerm$ & & \subdued{flat unification} \\
%         & $\mid$ & $R \, x_1 \ldots \, x_{k} $ & $ x_{j}\in V, x_j - distinct$ & \subdued{flat call} \\
% \end{tabular}
% \end{center}
% \end{frame}


% \begin{frame}[fragile]
%   \frametitle{Modes --- Data Flow}
% \begin{center}

% \begin{tabular}{rl}
%     \emph{Ground} term & \lstinline|S (S O)| \\
%   \emph{Free} variable & \lstinline|x|
% \end{tabular}

% \vfill

% Once a variable is \emph{ground}, it stays \emph{ground}
% \end{center}

% \vfill

% \begin{center}
% \[ \text{Mode} : \text{Inst} \mapsto \text{Inst} \] 
% \vfill

% \begin{tabular}{rl}
%   Mode \lstinline|I|: & \lstinline|ground| $\rightarrow$ \lstinline|ground| \\
%   Mode \lstinline|O|: & \lstinline|free| $\rightarrow$ \lstinline|ground|
% \end{tabular}
% \end{center}

% \end{frame}

% \begin{frame}[fragile]
%   \frametitle{Modded Unification Types}

% \begin{center}
% \begin{tabular}{rl}
%   assignment & $x^{\outmode} \equiv \Term^{\inmode} $ \\
%   guard      & $x^{\inmode}  \equiv \Term^{\inmode} $ \\
%   match      & $x^{\inmode}  \equiv \Term           $ \\
%   generator  & $x^{\outmode} \equiv \Term           $
% \end{tabular}
% \end{center}

% \end{frame}



% \begin{frame}[fragile]
%   \frametitle{Order in Conjunctions}
% \lstset{language=ocanren1,basicstyle={\ttfamily \footnotesize}}

\begin{center}
  \begin{tikzpicture}[
    every node/.style = {
        inner sep=0pt
      , shape=rectangle
      % , draw
      }
  ]
      % \draw[help lines] (-7.75,-3.5) grid (7.75,3.5);
      \node[inner sep=10pt, anchor=north west] (func) at (-7.75,3.5) {\begin{minipage}{0.35\columnwidth}
  \begin{lstlisting}[frame=tb]
 let rec mult$^o$ x y z = conde [
   (fresh (x$_1$ r$_1$)
     (x === S x$_1$) /\
     (add$^o$ y r$_1$ z) /\
     (mult$^o$ x$_1$ y r$_1$));
   ...]
\end{lstlisting}
\end{minipage}};

      \node[inner sep=10pt, anchor=north east] (slow) at (7.75,3.5) {\begin{minipage}{0.52\columnwidth}
  \begin{lstlisting}[frame=tb]
 multIIO$_1$ $::$ Nat -> Nat -> Stream Nat
 multIIO$_1$ (S x$_1$) y    = do
   (r$_1$, r) <- addIOO y
   multIII x$_1$ y r$_1$
   return r
 ...
 multIII $::$ Nat -> Nat -> Nat -> Stream ()
 multIII (S x$_1$) y z = do
   z$_1$ <- multIIO$_1$ x$_1$ y
   addIII y z$_1$ z
 multIII _ _ _ = Empty
 ...
\end{lstlisting}
\end{minipage}
};

      \node [draw=none, above=of slow.south, xshift=2.3cm] {\badBadge{$\Omega(x!)$}};

      \node[ellipse, scale=0.8,
          fill=badRed!75,
          draw=darkFastGreen,
          line width=0.4mm,
          inner sep=1pt,
          align=center, above=of slow.south, xshift=2.5cm, yshift=4cm]
        {\textbf{generate-and-test}};

      \node[inner sep=10pt, anchor=south west] (fast) at (-7.75,-3.5) {\begin{minipage}{0.42\columnwidth}
\begin{lstlisting}[frame=tb]
multIIO $::$ Nat -> Nat -> Stream Nat
multIIO (S x$_1$) y = do
  r$_1$ <- multIIO x$_1$ y
  addIIO y r$_1$
...
\end{lstlisting}
\end{minipage}
};

      \node [draw=none, anchor=south west] at ($(fast.south)+(0.5,0.8)$) {\goodBadge{$O(xy)$}};
      \node [draw=none, above=of fast.south, xshift=2.3cm] {\goodBadge{10x \linebreak faster}};
  \end{tikzpicture}
  \end{center}
  
% \end{frame}


% % \begin{frame}[fragile]
% %   \frametitle{Order in Conjunctions: Slow Version}
% % \begin{columns}
% %   \begin{column}[t]{0.4\textwidth}
% %     \begin{figure}[!t]
  \centering
  \begin{minipage}{\columnwidth}
    \begin{lstlisting}[label={mult},
                      %  caption={Multiplication relation},
                       captionpos=b,
                       frame=tb]
let rec mult$^o$ x y z = conde [
  ...
  (fresh (x' r')
    (x === S x') /\
    (add y r' z) /\
    (mult x' y r')
  )]
    \end{lstlisting}
  \end{minipage}
\end{figure}

% %   \end{column}
% %   \begin{column}[t]{0.55\textwidth}
% %     \begin{figure}[!t]
  \centering
  \begin{minipage}{\columnwidth}
    \begin{lstlisting}[label={mult_slow}, caption={Inefficient implementation of \lstinline{multo in in out} direciton}, captionpos=b, frame=tb]
multXY' :: Nat -> Nat -> Stream Nat
multXY' O     y     = return O
multXY' x     O     = return O
multXY' (S O) y     = return y
multXY' x     (S O) = return x
multXY' (S x') y    = do
  (r', r) <- addX y
  multXYZ x' y r'
  return r

multXYZ :: Nat -> Nat -> Nat -> Stream ()
multXYZ O      y     O = return ()
multXYZ x      O     O = return ()
multXYZ (S O)  y     z | y == z = return ()
multXYZ x      (S O) z | x == z = return ()
multXYZ (S x') y     z = do
  z' <- multXY' x' y
  addXYZ y z' z
multXYZ _ _ _ = Empty
    \end{lstlisting}
  \end{minipage}
\end{figure}

% %   \end{column}
% % \end{columns}

% % \begin{center}
% %   Premature grounding of \lstinline{z$_1$} leads to the \emph{generate-and-test} behavior

% %   $\Omega(x!)$ complexity.
% % \end{center}

% % \end{frame}


% % \begin{frame}[fragile]
% %   \frametitle{Order in Conjunctions: Faster Version}

% % \begin{columns}
% %   \begin{column}[t]{0.4\textwidth}
% %     \begin{figure}[!t]
  \centering
  \begin{minipage}{\columnwidth}
    \begin{lstlisting}[label={mult},
                      %  caption={Multiplication relation},
                       captionpos=b,
                       frame=tb]
let rec mult$^o$ x y z = conde [
  ...
  (fresh (x' r')
    (x === S x') /\
    (add y r' z) /\
    (mult x' y r')
  )]
    \end{lstlisting}
  \end{minipage}
\end{figure}

% %   \end{column}
% %   \begin{column}[t]{0.55\textwidth}
% %     \begin{figure}[!t]
  \centering
  \begin{minipage}{0.8\columnwidth}
    \begin{lstlisting}[frame=tb]
 multIIO $::$ Nat -> Nat -> Stream Nat
 multIIO (S x$_1$) y = do
   r$_1$ <- multIIO x$_1$ y
   addIIO y r$_1$
 ...
    \end{lstlisting}
  \end{minipage}
\end{figure}

% %   \end{column}
% % \end{columns}
% % \begin{center}
% %     $O(xy)$ complexity, $10$x faster than relational version
% % \end{center}
% % \end{frame}

% \begin{frame}[fragile]
%   \frametitle{Mode Inference: Ordering Heuristic}
% \begin{center}
% Priority:
% \end{center}

% \vfill

% \begin{center}
%   \begin{minipage}{0.4\textwidth}
%     \begin{enumerate}
%       \item Guard
%       \item Assignment
%       \item Match
%       \item Recursion, same direction
%       \item Call, some args ground
%       \item Unification-generator
%       \item Call, all args free
%     \end{enumerate}
%   \end{minipage}
% \end{center}


% \end{frame}

% \begin{frame}[fragile]
%   \frametitle{Ordering Heuristic: Example}
% \begin{center}
  \begin{tikzpicture}[
    every node/.style = {
        inner sep=0pt
      , shape=rectangle
      % , draw
      }
  ]
      % \draw[help lines] (-7.75,-3.5) grid (7.75,3.5);
      \node[inner sep=10pt, anchor=north] (func) at (-7.75,3.5) {\begin{minipage}{0.35\columnwidth}
  \begin{lstlisting}[frame=tb]
 let rec mult$^o$ x y z = conde [
   (fresh (x$_1$ r$_1$)
     (x === S x$_1$) /\
     (add$^o$ y r$_1$ z) /\
     (mult$^o$ x$_1$ y r$_1$));
   ...]
\end{lstlisting}
\end{minipage}};

      \node[inner sep=10pt, anchor=south] (fast) at (-7.75,-3.5) {\begin{minipage}{0.42\columnwidth}
\begin{lstlisting}[frame=tb]
multIIO $::$ Nat -> Nat -> Stream Nat
multIIO (S x$_1$) y = do
  r$_1$ <- multIIO x$_1$ y
  addIIO y r$_1$
...
\end{lstlisting}
\end{minipage}
};

  \end{tikzpicture}
  \end{center}
  

% \end{frame}

% \begin{frame}[fragile]
%   \frametitle{Functional Conversion: Intermediate Language}
% \begin{center}

% \begin{tabular}{lcl @{\hspace{3\tabcolsep}} l }
%     $\Fun$ & $=$ & $\Fun <|> \dots <|> \Fun$ & \subdued{interleaving}\\
%                & $\mid$ & $\LIST{\left(\overline{V} \leftarrow \Fun\right)} $ & \subdued{monadic bind of streams}\\
%                & $\mid$ & $\Rtrn \LIST{\Term}$ & \subdued{return a tuple of terms}\\
%                & $\mid$ & $V == \Term$ & \subdued{equality check}\\
%                & $\mid$ &  $case \ V \ of \ \Term \rightarrow \Fun$& \subdued{match a variable against a pattern}\\
%                & $\mid$ & $R \ \overline{V} \ \overline{Gen_{G}}$ & \subdued{function call}\\
%                & $\mid$ & $\Gen_{G}$ & \subdued{generator}
% \end{tabular}
% \end{center}
% \end{frame}

% \begin{frame}[fragile]
%   \frametitle{Functional Conversion into Intermediate Language}
% \begin{center}
% \begin{tabular}{rcl}
%   Disjunction   & $\rightarrow$ & $\Fun \Sum \ldots \Sum \Fun$ \\ && \\
%   Conjunction   & $\rightarrow$ & $ \LIST{\left(\overline{V} \leftarrow \Fun\right)}$ \\ && \\
%   Relation call & $\rightarrow$ & $ R\ \overline{V}\ \overline{Gen_{G}}$ \\ && \\
%   Unification   & $\rightarrow$ & $\Rtrn \LIST{\Term}$ \\
%                 & $|$           & $case\ V\ of\ \Term \to \Fun$ \\
%                 & $|$           & $V \equiv \Term$ \\
%                 & $|$           & $\Gen_{G}$
% \end{tabular}
% \end{center}
% \end{frame}


% \begin{frame}[fragile]
%   \frametitle{Functional Conversion: Generators}
% % \begin{center}
% %   In the untyped miniKanren it is only possible to generate \emph{all terms}
% % \end{center}

% % \vfill

% % \begin{center}
% %   Instead pass generators to functions as additional arguments


% %   It is up to the user what generator to pass
% % \end{center}

% % \begin{figure}[!t]
  \centering
  \begin{minipage}{0.59\columnwidth}
    \begin{lstlisting}[frame=tb]
 addIOO $::$ Nat -> Stream Nat -> Stream (Nat, Nat)
 addIOO x gen$_z$ =
   case x of
     O -> do
       z <- gen$_z$
       return (z, z)
     S x$_1$ -> do
       (y, z$_1$) <- addIOO x$_1$ gen$_z$
       return (y, S z$_1$)
    \end{lstlisting}
  \end{minipage}
\end{figure}

% \begin{center}
  \begin{tikzpicture}[
    every node/.style = {
      shape=rectangle,
      rounded corners=0.05cm,
      draw,
      align=center,
      minimum size=5mm},
    node distance=1.3em,
    anchor=center
  ]
    \node[inner sep=10pt,draw=none] (code) at (0,0) {
      \begin{minipage}{0.59\columnwidth}
      \begin{lstlisting}[frame=tb]
 addIOO $::$ Nat -> Stream Nat -> Stream (Nat, Nat)
 addIOO x gen$_y$ =
   case x of
     O -> do
       y <- gen$_y$
       return (y, y)
     S x$_1$ -> do
       (y, z$_1$) <- addIOO x$_1$ gen$_y$
       return (y, S z$_1$)
       \end{lstlisting}
      \end{minipage}
    };
    \draw [teal,-Stealth] ($(code.center)+(-1.6,0.2)$) to [out=40,in=-30] node[below, draw=none, teal] {} ($(code.north)+(-2,-1.1)$);
    \draw [teal,-Stealth] ($(code.center)+(-1.6,0.2)$) to [out=40,in=130] node[below, draw=none, teal] {} ($(code.south)+(0.5,1.4)$);
      
  \end{tikzpicture}
  \end{center}
  
% \[ add^{\circ}\ 1\ y\ z = \left[\{z \to S\ y\} \right] \] 
% \end{frame}

% \begin{frame}[fragile]
%   \frametitle{Functional Conversion: Generators}
% % \begin{center}
% % We pass a generator for every variable in \emph{rhs} of a unification-generator
% % \end{center}

% % \begin{center}
% % Generators used in calls should be passed to the parent function
% % \end{center}

% % \begin{center}
% % In a typed version, it should be possible to automatically derive generators from types
% % \end{center}

% % \begin{figure}[!t]
  \centering
  \begin{minipage}{\columnwidth}
    \begin{lstlisting}[label={mult_x_gen},
                      %  caption={Efficient implementation of \lstinline{multo in in out} direciton},
                       captionpos=b,
                       frame=tb]
multOIO :: Nat -> Stream Nat -> Stream Nat
multOIO y gen_add$_z$ =
  return (O, O) `mplus`
  do
    (z$_1$, z) <- addIOO y gen_add$_z$
    x <- muloOII y z$_1$
    return (S x, z)
    \end{lstlisting}
  \end{minipage}
\end{figure}

% \begin{center}
  \begin{tikzpicture}[
    every node/.style = {
      shape=rectangle,
      rounded corners=0.05cm,
      draw,
      align=center,
      minimum size=5mm},
    node distance=1.3em,
    anchor=center
  ]
    \node[inner sep=10pt,draw=none] (code) at (0,0) {
      \begin{minipage}{0.59\columnwidth}
      \begin{lstlisting}[frame=tb]
 multOIO $::$ Nat -> Stream Nat -> Stream Nat
 multOIO y gen_add$_y$ =
   return (O, O) <$\mid$>
   do
     (z$_1$, z) <- addIOO y gen_add$_y$
     x <- muloOII y z$_1$
     return (S x, z)
       \end{lstlisting}
      \end{minipage}
    };
    \draw [teal,-Stealth] ($(code.center)+(0.3,-0.1)$) to [out=130,in=-25] node[below, draw=none, teal] {} ($(code.north)+(-1.1,-1.1)$);
      
  \end{tikzpicture}
  \end{center}
  
% \end{frame}

% % \begin{frame}[fragile]
% %   \frametitle{Functional Conversion into the Target Languages}
% % \begin{columns}
% %   \begin{column}[t]{0.45\textwidth}

% % \begin{center}
% %   \haskell
% % \end{center}

% % \begin{center}
% %   TemplateHaskell to generate code
% % \end{center}

% % \begin{center}
% %   Stream monad
% % \end{center}

% % \begin{center}
% %   do-notation
% % \end{center}

% % \begin{center}

% % \end{center}

% %   \end{column}
% %   \begin{column}[t]{0.45\textwidth}
% % \begin{center}
% %   \ocaml
% % \end{center}
% % \begin{center}
% %   Hand-crafted (not so) pretty-printer
% % \end{center}

% % \begin{center}
% %   Stream monad
% % \end{center}

% % \begin{center}
% %   let*
% % \end{center}

% % \begin{center}
% %   Taking extra care to ensure laziness
% % \end{center}

% %   \end{column}
% % \end{columns}
% % \end{frame}

% \begin{frame}[fragile]
%   \frametitle{Relational Sort}
% \begin{columns}
%   \begin{column}[t]{0.49\textwidth}
%     
\lstset{moredelim=[is][\bfseries]{[*}{*]}}

\begin{figure}[h]
  \centering
  \begin{minipage}{0.65\columnwidth}
    \begin{lstlisting}[frame=tb,language=ocanren1]
 let rec sort$^o$ x y =
   (x === [] /\ y === []) \/ 
   (fresh (s xs xs$_1$)
     y === s :: xs$_1$ /\
     smallest$^o$ x s xs /\
     sort$^o$ xs xs$_1$)
    \end{lstlisting}
  \end{minipage}
\end{figure}


%     \vfill

%     \begin{center}
%       \begin{minipage}{0.4\textwidth}
%         \begin{itemize}
%           \item[\happyCheck] sorting
%           \item[\timeout] permutations
%         \end{itemize}
%       \end{minipage}
%     \end{center}

%   \end{column}
%   \begin{column}[t]{0.49\textwidth}
%     \begin{figure}[h]
  \centering
  \begin{minipage}{0.87\columnwidth}
    \begin{lstlisting}[frame=tb]
 let rec sort$^o$ x y = conde [
    (x === [] /\ y === []);
    (fresh (s xs xs$_1$)
      y === s :: xs$_1$ /\
      [*sort$^o$ xs xs$_1$ /\
      smallest$^o$ x s xs)*]]
    \end{lstlisting}
  \end{minipage}
\end{figure}


%     \vfill

%     \begin{center}
%       \begin{minipage}{0.4\textwidth}
%         \begin{itemize}
%           \item[\timeout] sorting
%           \item[\happyCheck] permutations
%         \end{itemize}
%       \end{minipage}
%     \end{center}


%   \end{column}
% \end{columns}
% \end{frame}

% \begin{frame}[fragile]
%   \frametitle{Relational Sort: Sorting}
%     \begin{table}[]
\begin{tabular}{l||c|r||r}
                & \multicolumn{2}{c||}{Relation} & \multicolumn{1}{c}{\multirow[t]{2}{*}{Function}} \\ \cline{2-3}
                & \multicolumn{1}{l|}{\begin{tabular}[c]{@{}l@{}}\texttt{sorto}\\ \texttt{smallesto}\end{tabular}} & \multicolumn{1}{l||}{\begin{tabular}[c]{@{}l@{}}\texttt{smallesto}\\ \text{sorto}\end{tabular}} & \multicolumn{1}{c}{}                            \\ \hline
\texttt{{[}3;2;1;0{]}}   & \multicolumn{1}{r|}{0.077s}                                                    & 0.004s                                                                         & 0.000s                                          \\
\texttt{{[}4;3;2;1;0{]}} & \timeout timeout                                                                       & 0.005s                                                                         & 0.000s                                          \\
\texttt{{[}31;...;0{]}}  & \timeout timeout                                                                       & 1.058s                                                                         & 0.006s                                          \\
\texttt{{[}262;...;0{]}} & \timeout timeout                                                                       & \multicolumn{1}{c||}{\timeout timeout}                                                   & 1.045s
\end{tabular}
\end{table}


% \end{frame}


% \begin{frame}[fragile]
%   \frametitle{Relational Sort: Generating Permutations}
%     \begin{table}[]
\begin{tabular}{l||c|r||r}
                & \multicolumn{2}{c||}{Relation} & \multicolumn{1}{c}{\multirow[t]{2}{*}{Function}} \\ \cline{2-3}
                & \multicolumn{1}{l|}{\begin{tabular}[c]{@{}l@{}}\texttt{smallesto}\\ \texttt{sorto}\end{tabular}} & \multicolumn{1}{l||}{\begin{tabular}[c]{@{}l@{}}\texttt{sorto}\\ \text{smallesto}\end{tabular}} & \multicolumn{1}{c}{}                            \\ \hline
\texttt{{[}0;1;2{]}}   & \multicolumn{1}{r|}{0.013s}                                                    & 0.004s                                                                         & 0.004s                                          \\
\texttt{{[}0;1;2;3{]}} & \timeout timeout                                                                       & 0.005s                                                                         & 0.005s                                          \\
\texttt{{[}0;...;6{]}}  & \timeout timeout                                                                       & 0.999s                                                                         & 0.021s                                          \\
\texttt{{[}0;...;8{]}} & \timeout timeout                                                                       & \multicolumn{1}{c||}{\timeout timeout}                                                   & 1.543s
\end{tabular}
\end{table}

% \end{frame}

% \begin{frame}[fragile]
%     \frametitle{Possible speedups}

%     \begin{center}
%         \begin{minipage}{0.4\textwidth}
%             \begin{itemize}
%                 \item Determinism
%                 \item Partial deduction
%             \end{itemize}
%         \end{minipage}
%     \end{center}
% \end{frame}

% \begin{frame}[fragile]
%   \frametitle{\lstinline[basicstyle=\Large]{Maybe} for Semi-Determinism}
% \begin{center}
%   \begin{minipage}{0.43\textwidth}
%     \begin{center}
  \begin{tikzpicture}[
    every node/.style = {
      shape=rectangle,
      rounded corners=0.05cm,
      draw,
      align=center,
      minimum size=5mm},
    node distance=1.3em,
    anchor=center
  ]
    \node[inner sep=10pt,draw=none] (code) at (0,0) {
      \begin{minipage}{\columnwidth}
      \begin{lstlisting}[]

 ~\texttt{muloOII $::$ Nat $\to$ Nat $\to$ Stream Nat}~
 muloOII x1 x2 =
     zero <$\mid$> positive
   where
     zero = do
       guard (x2 == O)
       return O
     positive = do
       x4 <- addoIOI x1 x2
       S <$\$$> muloOII x1 x4
       \end{lstlisting}
      \end{minipage}
    };
  \end{tikzpicture}
  \end{center}
  
%     % \begin{figure}[!t]
  \centering
  \begin{minipage}{\columnwidth}
    \begin{lstlisting}[frame=tb]
muloOII :: Nat -> Nat -> Stream Nat
muloOII x1 x2 = zero `mplus` positive
  where
    zero = do
      guard (x2 == O)
      return O
    positive = do
      x4 <- addoIOI x1 x2
      S <$\$$> muloOII x1 x4
    \end{lstlisting}
  \end{minipage}
\end{figure}
%   \end{minipage}
% \end{center}
% \end{frame}


% \begin{frame}[noframenumbering]
%   \frametitle{\lstinline[basicstyle=\Large]{Maybe} for Semi-Determinism}
%   \begin{center}
%   \begin{minipage}{0.43\textwidth}
%     \begin{center}
  \begin{tikzpicture}[
    every node/.style = {
      shape=rectangle,
      rounded corners=0.05cm,
      draw,
      align=center,
      minimum size=5mm},
    node distance=1.3em,
    anchor=center
  ]
    \node[inner sep=10pt,draw=none] (code) at (0,0) {
      \begin{minipage}{\columnwidth}
      \begin{lstlisting}[]
 ~\texttt{muloIIO $::$ Nat $\to$ Nat $\to$ Maybe Nat}~
 ~\lststrikethrough{muloIIO $::$ Nat $\to$ Nat $\to$ Stream Nat}~
 muloIIO x y =
     zero <$\mid$> positive
   where
     zero = do
       guard (x == O)
       return O
     positive = do
       (S $x_1$) <- return x
       $z_1$ <- muloIIO $x_1$ y
       z <- addoIIO $z_1$ y
       return z
       \end{lstlisting}
      \end{minipage}
    };
    \node [draw=none, above=of code.south, xshift=2.5cm, yshift=2.5cm] {\goodBadge{10x \linebreak faster}};
  \end{tikzpicture}
  \end{center}
  

%     % \lstset{moredelim=[is][\sout]{|}{|}}

\begin{figure}[!t]
  \centering
  \begin{minipage}{\columnwidth}
    \begin{lstlisting}[]
 ~\texttt{muloOII $::$ Nat $\to$ Nat $\to$ Maybe Nat}~
 ~\lststrikethrough{muloOII $::$ Nat $\to$ Nat $\to$ Stream Nat}~
 muloOII x1 x2 =
     zero <$\mid$> positive
   where
     zero = do
       guard (x2 == O)
       return O
     positive = do
       x4 <- addoIOI x1 x2
       S <$\$$> muloOII x1 x4
    \end{lstlisting}
  \end{minipage}
\end{figure}

%   \end{minipage}
% \end{center}
% \end{frame}

% \begin{frame}[fragile]
%   \frametitle{Need for Partial Deduction}

% \begin{center}
% \begin{minipage}{0.3\textwidth}
%   \lstinline{run q (eval$^o$ q true)}

%   \vspace{0.5cm}

%   \lstinline{run q (eval_true$^o$ q)}
% \end{minipage}


% \end{center}


% \end{frame}

% \begin{frame}[fragile]
%   \frametitle{Conclusion}
% Conclusion
%   \begin{itemize}
%     \item We presented a functional conversion scheme
%     \item The conversion speeds up implementations considerably
%     \item We implemented the conversion scheme in Haskell
%   \end{itemize}

% \vfill

% We are currently working on
%   \begin{itemize}
%     \item Determinism check
%     \item Integration with partial deduction
%     \item Integration into the framework of using relational interpreters for solving
%   \end{itemize}
% \end{frame}





% \appendix

% \begin{frame}[fragile]
%     \frametitle{Functional Conversion: Example}
%     \begin{columns}
%         \begin{column}[t]{0.45\textwidth}
%             \begin{figure}[!t]
  \centering
  \begin{minipage}{0.7\columnwidth}
    \begin{lstlisting}[frame=tb]
 let rec add$^o$ x y z = conde [
   (x === O /\ y === z);
   (fresh (x$_1$ z$_1$)
     (x === S x$_1$ /\
      add$^o$ x$_1$ y z$_1$ /\
      z === S z$_1$) ) ]
    \end{lstlisting}
  \end{minipage}
\end{figure}
%         \end{column}
%         \begin{column}[t]{0.55\textwidth}
%             \begin{figure}[h]
\centering
\begin{minipage}{0.9\columnwidth}
  \begin{lstlisting}[frame=tb]
data Term = O | S Term
addoIIO :: Term -> Term -> Stream Term
addoIIO x y = msum [
    do {
        guard (x == O);
        z <- return y;
        return z
    },
    do {
        S $x_1$ <- return x;
        $z_1$ <- addoIIO $x_1$ y;
        z <- return (S $z_1$);
        return z
    }
]
  \end{lstlisting}
\end{minipage}
\end{figure}


%         \end{column}
%     \end{columns}
% \end{frame}

\backupbegin

\backupend

\end{document}
